\documentclass[10pt]{article}
\usepackage[a4paper,margin=18mm]{geometry}
\usepackage[no-math]{fontspec}
\setmainfont{Latin Modern Roman}
\usepackage{amsmath,amssymb,amsthm,mathrsfs,enumitem}
\usepackage[all,cmtip]{xy}
\usepackage{hyperref}
\hypersetup{hidelinks}
\newtheorem{theorem}{Theorem}[section]
\newtheorem{lemma}[theorem]{Lemma}
\newtheorem{proposition}[theorem]{Proposition}
\newtheorem{corollary}[theorem]{Corollary}
\theoremstyle{definition}
\newtheorem{definition}[theorem]{Definition}
\newtheorem{remark}[theorem]{Remark}
\newtheorem{example}[theorem]{Example}
\newcounter{intro}
\renewcommand{\theintro}{\Alph{intro}}
\newtheorem{intro-theorem}[intro]{Theorem}
\newtheorem{intro-corollary}[intro]{Corollary}
\renewcommand{\thefootnote}{(\arabic{footnote})}
\renewcommand{\qedsymbol}{\ensuremath{\square}}
\numberwithin{equation}{section}
\emergencystretch=3em
\pdfstringdefDisableCommands{\def\jepPubBold#1{#1}\def\jepPubScript#1{#1}}
\DeclareRobustCommand{\jepPubScript}[1]{\begingroup\def\jepArg{#1}\def\jepTestD{D}\def\jepTestH{H}\def\jepTestI{I}\def\jepTestM{M}\def\jepTestN{N}\def\jepTestO{O}\def\jepTestP{P}\ifx\jepArg\jepTestD\mathscr{D}\else\ifx\jepArg\jepTestH\mathscr{H}\else\ifx\jepArg\jepTestI\mathscr{I}\else\ifx\jepArg\jepTestM\mathscr{M}\else\ifx\jepArg\jepTestN\mathscr{N}\else\ifx\jepArg\jepTestO\mathscr{O}\else\ifx\jepArg\jepTestP\mathscr{P}\else\PackageError{scoped-published-font}{Unsupported literal in jepPubScript}{Only inventoried published arguments are allowed.}\fi\fi\fi\fi\fi\fi\fi\endgroup}
\DeclareRobustCommand{\jepPubBold}[1]{\begingroup\def\jepArg{#1}\def\jepTestA{A}\def\jepTestC{C}\def\jepTestL{L}\def\jepTestQ{Q}\def\jepTestR{R}\def\jepTestZ{Z}\ifx\jepArg\jepTestA\mathbf{A}\else\ifx\jepArg\jepTestC\mathbf{C}\else\ifx\jepArg\jepTestL\mathbf{L}\else\ifx\jepArg\jepTestQ\mathbf{Q}\else\ifx\jepArg\jepTestR\mathbf{R}\else\ifx\jepArg\jepTestZ\mathbf{Z}\else\PackageError{scoped-published-font}{Unsupported literal in jepPubBold}{Only inventoried published arguments are allowed.}\fi\fi\fi\fi\fi\fi\endgroup}
\newcommand{\jepPubBullet}{{\scriptscriptstyle\bullet}}
\begin{document}
\begin{center}{\Large Stable item 1906: checked published transcription}\end{center}
\paragraph{Source and attribution.}
Mircea Mustaţǎ and Mihnea Popa, \emph{Hodge ideals for $\mathbf Q$-divisors: birational approach}, Journal de l’École polytechnique -- Mathématiques 6 (2019), 283--328, DOI 10.5802/jep.94. \href{https://jep.centre-mersenne.org/articles/10.5802/jep.94/}{Publisher article}; \href{https://creativecommons.org/licenses/by/4.0/}{CC BY 4.0}.
\paragraph{Selected independent problem.}
Original Theorem 8.1, for a smooth complex variety, a nonzero regular function $h$, positive rational $\alpha$, $D=\alpha\operatorname{div}(h)$, and a log resolution isomorphic over the complement of its support. All original introductory material and Sections 1--7, the complete Section 8 setup, Theorem 8.1 and proof, adjacent Remark 8.2, and all 19 bibliography entries are retained. The original full 47-page publisher PDF is retained separately. Source Sections 9--15 and the local Nakano corollary are outside this selected unit. Supporting results are uncounted context.
\paragraph{Difficulty (editorial): H3.}
Equivariant cyclic covers define the filtered rational-power $\mathscr D$-module. The twisted logarithmic Spencer complex is proved exact by an explicit local normal-form induction. Toric normalization then determines the round-up lattice and eigenspace filtration, before strictness yields the birational vanishing and exact filtered image computing every Hodge ideal.
\paragraph{Representation disclosure.}
This is an editable publisher-authority transcription, not native author TeX. The nonexclusive arXiv v2 archive was read only as an unexecuted private aid. No private archive, preamble, class, comment, or unpublished version difference is exported or loaded. Original Section 2.11 prints a positive exponent in the auxiliary cover, Section 2.12 prints $gh^m$ in its displayed map, Section 6 prints $h^\alpha$ in one prose filtration reference, and the proof of Theorem 8.1 closes with ``proposition''; these source literals are preserved, without mathematical repair.
\clearpage
\begin{center}{\Large Hodge ideals for $\jepPubBold{Q}$-divisors: birational approach}\\[1ex]Mircea Mustaţǎ \& Mihnea Popa\end{center}
\paragraph{Abstract.}
We develop the theory of Hodge ideals for $\jepPubBold{Q}$-divisors by means of log resolutions, extending our previous work on reduced hypersurfaces. We prove local (non-)triviality criteria and a global vanishing theorem, as well as other analogues of standard results from the theory of multiplier ideals, and we derive a new local vanishing theorem. The connection with the $V$-filtration is analyzed in a sequel.
\paragraph{Résumé (Idéaux de Hodge pour des $\jepPubBold{Q}$-diviseurs : approche birationnelle).}
Nous développons la théorie des idéaux de Hodge pour les $\jepPubBold{Q}$-diviseurs à l’aide de log résolutions, généralisant notre précédent travail sur les hypersurfaces réduites. Nous obtenons des critères de (non) trivialité locale et un théorème d’annulation global, ainsi que d’autres analogues de résultats standard de la théorie des idéaux multiplicateurs, et nous en déduisons un nouveau théorème d’annulation local. Nous analysons la relation avec la $V$-filtration dans un autre article.
\section*{A. Introduction}
In this paper we continue the study of Hodge ideals initiated in \cite{MP1}, \cite{MP2}, by considering an analogous theory 
for arbitrary $\jepPubBold{Q}$-divisors. The emphasis here is on a birational definition and study of Hodge ideals, while the companion 
paper \cite{MP3} is devoted to a study based on their connection with the $V$-filtration, inspired by \cite{Saito-MLCT}. Both 
approaches turn out to provide crucial information towards a complete understanding of these objects.

Let $X$ be a smooth complex variety. If $D$ is reduced divisor on $X$, the Hodge ideals $I_k (D)$, with $k \geqslant 0$, 
are defined in terms of the Hodge filtration 
on the $\jepPubScript{D}_X$-module $\jepPubScript{O}_X(*D)$ of functions with poles of arbitrary order along $D$. Indeed, this 
$\jepPubScript{D}_X$-module
underlies a mixed Hodge module on $X$, and therefore comes with a Hodge filtration $F_\jepPubBullet \jepPubScript{O}_X (*D)$, which satisfies
$$F_k \jepPubScript{O}_X(*D) = I_k (D) \otimes \jepPubScript{O}_X\big( (k+1)D\big)\quad\text{for all}\,\,\,\, k \geqslant 0.$$
See \cite{MP1} for details, and for an extensive study of the ideals $I_k (D)$.

Our goal here is to provide a similar construction and study in the general case. A natural device for dealing with the 
fact that fractional divisors are not directly related to Hodge theory is to use new objects derived from covering constructions.
Let $D$ be an arbitrary effective $\jepPubBold{Q}$-divisor on $X$. Locally, we can write $D = \alpha H$, for 
some $\alpha \in \jepPubBold{Q}_{> 0}$ and $H = {\rm div} (h)$, the divisor of a nonzero regular function; we also denote by $Z$ the support of $D$.
A well-known construction associates to this data a twisted version of the localization $\jepPubScript{D}$-module above, namely
$$\jepPubScript{M} (h^{-\alpha}) : = \jepPubScript{O}_X (*Z) h^{-\alpha},$$
that is the rank $1$ free $\jepPubScript{O}_X(*Z)$-module 
with generator the symbol $h^{-\alpha}$, on which a derivation $D$ of $\jepPubScript{O}_X$ acts by
$$D(wh^{-\alpha}) :=\Big(D(w) - \alpha w\frac{ D(h)}{h}\Big)h^{-\alpha}.$$
It turns out that this $\jepPubScript{D}_X$-module can be endowed with a natural filtration $F_k \jepPubScript{M} (h^{-\alpha})$, with 
$k \geqslant 0$, which makes it a filtered direct summand of a $\jepPubScript{D}$-module underlying a mixed Hodge module on $X$; see 
Section~\ref{coverings}. This plays a role 
analogous to the Hodge filtration, and just as in the reduced case one can show that $F_k \jepPubScript{M} (h^{-\alpha}) \subseteq 
\jepPubScript{O}_X (kZ ) h^{-\alpha}$. This is done in Section~\ref{smooth} and Section~\ref{definition}, by first analyzing the case when $Z$ is a smooth divisor 
(in this case, if $\lceil D\rceil=Z$, then
the inclusion is in fact an equality). It is therefore natural to define the \emph{$k$-th Hodge ideal} of $D$ by the formula
$$F_k \jepPubScript{M} (h^{-\alpha}) = I_k (D) \otimes_{\jepPubScript{O}_X} \jepPubScript{O}_X (kZ ) h^{-\alpha}.$$ 

Similarly to \cite{MP1}, one of our main goals here is to study Hodge ideals of $\jepPubBold{Q}$-divisors by means of log resolutions. 
To this end, let $f\colon Y\to X$ be a log resolution of the pair $(X, D)$ that is an isomorphism over $U=X\smallsetminus Z$,  
and denote $g=h\circ f$. There is a filtered isomorphism
$$
\big(\jepPubScript{M}(h^{-\alpha}), F \big)\simeq f_+\big(\jepPubScript{M}(g^{-\alpha}), F\big).
$$
Denoting $G=f^*D$ and $E = {\rm Supp}(G)$, so that $E$ is a simple normal crossing divisor, 
it turns out that there exists a complex on $Y$:
$$C^{\jepPubBullet}_{g^{-\alpha}} (- \lceil G\rceil):\,\,0\longrightarrow\jepPubScript{O}_Y(-\lceil G\rceil)\otimes_{\jepPubScript{O}_Y} \jepPubScript{D}_Y\longrightarrow\jepPubScript{O}_Y(-\lceil G\rceil)\otimes_{\jepPubScript{O}_Y}\Omega^1_Y(\log E)\otimes_{\jepPubScript{O}_Y}\jepPubScript{D}_Y$$
$$\longrightarrow\cdots\longrightarrow\jepPubScript{O}_Y(-\lceil G\rceil)\otimes_{\jepPubScript{O}_Y}\omega_Y(E)\otimes_{\jepPubScript{O}_Y}\jepPubScript{D}_Y\longrightarrow 0,$$
which is placed in degrees $-n,\ldots,0$, whose differential is described in Section~\ref{complex_for_SNC}.
This complex has a natural filtration given, for $k\geqslant 0$, by subcomplexes 
$$F_{k-n}C^{\jepPubBullet}_{g^{-\alpha}} (- \lceil G\rceil): = 
0 \rightarrow \jepPubScript{O}_Y(-\lceil G\rceil)\otimes  F_{k-n} \jepPubScript{D}_Y$$
$$\rightarrow  \jepPubScript{O}_Y(-\lceil G\rceil)\otimes\Omega_Y^1(\log E) \otimes  F_{k-n+1} \jepPubScript{D}_Y \rightarrow 
\cdots \rightarrow  \jepPubScript{O}_Y(-\lceil G\rceil)\otimes \omega_Y(E) \otimes  F_k \jepPubScript{D}_Y\rightarrow 0.$$

\noindent
Extending \cite[Prop. 3.1]{MP1}, we show in Proposition \ref{prop_SNC_complex} and Proposition \ref{Hodge_ideals_SNC} 
that there is a filtered quasi-isomorphism
$$
\big(C^{\jepPubBullet}_{g^{-\alpha}} (- \lceil G\rceil), F \big) \simeq \big(\jepPubScript{M}_r (g^{-\alpha}), F \big),
$$
where $\jepPubScript{M}_r (g^{-\alpha})$ is the filtered right $\jepPubScript{D}_Y$-module associated to $\jepPubScript{M} (g^{-\alpha})$.
Thus one can use $\big(C^{\jepPubBullet}_{g^{-\alpha}} (- \lceil G\rceil), F\big)$ as a concrete representative for computing the filtered 
$\jepPubScript{D}$-module pushforward of $\big(\jepPubScript{M}_r (g^{-\alpha}), F\big)$, hence for computing 
the ideals $I_k (D)$. More precisely, we have 
$$R^0f_*F_{k-n}\big(C^{\jepPubBullet}_{g^{-\alpha}} (- \lceil G\rceil)\otimes_{\jepPubScript{D}_Y}\jepPubScript{D}_{Y\to X}\big)\simeq
h^{-\alpha}\omega_X (kZ )\otimes_{\jepPubScript{O}_X}I_k(D).$$
See Theorem \ref{formula_log_resolution} for a complete picture regarding this push-forward operation.

This fact, together with special properties of the filtration on $\jepPubScript{D}$-modules underlying mixed Hodge modules, leads to 
our main results on Hodge ideals, which are collected in the following:

\begin{intro-theorem}\label{main_list}
In the set-up above, the Hodge ideals $I_k (D)$ satisfy:


\noindent
(i) $I_0 (D)$ is the multiplier ideal $\jepPubScript{I} \big((1 - \varepsilon)D\big)$, 
so in particular $I_0 (D) = \jepPubScript{O}_X$ if and only if the pair $(X, D)$ is log canonical; see Section~9.


\noindent
(ii) If $Z$ has simple normal crossings, then 
$$I_k (D) = I_k (Z) \otimes \jepPubScript{O}_X (Z - \lceil D \rceil),$$
while $I_k (Z)$ can be computed explicitly as in \cite[Prop.~8.2]{MP1}; see Section~\ref{SNC}.  In particular, if 
$Z$ is smooth, then $I_k (D) = \jepPubScript{O}_X (Z - \lceil D \rceil)$ for all $k$; cf. also Corollary 11.12.



\noindent
(iii) The Hodge filtration is generated at level $n-1$, where $n = \dim X$, i.e. 
$$F_\ell \jepPubScript{D}_X\cdot\big(I_k(D)\otimes \jepPubScript{O}_X(kZ)h^{-\alpha}\big) = I_{k+\ell}(D)\otimes \jepPubScript{O}_X \big((k+\ell)Z\big)h^{-\alpha}$$
for all $k \geqslant n-1$ and $\ell\geqslant 0$; see Section~10.


\noindent
(iv) There are non-triviality criteria for $I_k (D)$ at a point $x \in D$ in terms of the multiplicity of $D$ at $x$; 
see Section~11.



\noindent
(v) If $X$ is projective, $I_k(D)$ satisfy a vanishing theorem analogous to Nadel Vanishing for multiplier ideals; 
see Section~12.


\noindent
(vi) If $Y$ is a smooth divisor in $X$ such that $Z|_Y$ is reduced, then $I_k (D)$ satisfy 
$$I_k (D|_Y) \subseteq I_k(D) \cdot \jepPubScript{O}_Y,$$
with equality when $Y$ is general; see Section~13 for a more general statement.


\noindent
(vii) If $ X \to T$ is a smooth family with a section $s\colon T \to X$, and $D$ is a relative divisor on $X$ that satisfies a suitable condition
(see Section~14 for the precise statement) then 
$$\{ t \in T ~|~ I_k (D_t) \not\subseteq \mathfrak{m}_{s(t)}^q\}$$
is an open subset of $T$, for each $q \geqslant 1$.



\noindent
(viii) If $D_1$ and $D_2$ are $\jepPubBold{Q}$-divisors with supports $Z_1$ and $Z_2$, such that $Z_1 +  Z_2$ is also reduced, then  the subadditivity property
$$I_k (D_1 + D_2) \subseteq I_k (D_1)\cdot I_k (D_2)$$
holds; see Section~15 for a more general statement.
\end{intro-theorem}


For comparison, the list of properties of Hodge ideals in the case when $D$ is reduced is summarized in \cite[\S4]{Popa}. While much of the story carries over to the setting of $\jepPubBold{Q}$-divisors -- besides of course the connection with the classical Hodge theory of the complement $U = X \smallsetminus D$, which only makes sense in the reduced case -- there are a few significant points where the picture becomes more intricate. For instance, the bounds for the generation level of the Hodge filtration can become worse.  
Moreover, we do not know 
whether the inclusions $I_k (D) \subseteq I_{k-1} (D)$ continue to hold for arbitrary $\jepPubBold{Q}$-divisors. New phenomena appear as well: unlike in the case of multiplier ideals, for 
rational numbers $\alpha_1 < \alpha_2$, usually the ideals $I_k (\alpha_1 Z)$ and $I_k (\alpha_2 Z)$ cannot be 
compared for $k \geqslant 1$; see for instance Example 10.5.

It turns out however that most of these issues disappear if one works modulo the ideal of the hypersurface, at least for 
rational multiples of a reduced divisor.
This, as well as other basic facts, is addressed in the sequel \cite{MP3}, which 
studies Hodge ideals from a somewhat different point of view, namely by comparing them to
the (microlocal) $V$-filtration induced on $\jepPubScript{O}_X$ by $h$. This is inspired by the work of Saito \cite{Saito-MLCT} in the reduced case. In the statement below we summarize some of these properties, which complement the results in Theorem \ref{main_list}, but which we do not know how to obtain with the methods of this paper.


\begin{intro-theorem}[\cite{MP3}]\label{other_paper}
Let $D = \alpha Z$, where $Z$ is a reduced divisor and $\alpha \in \jepPubBold{Q}_{> 0}$. Then the following hold:
\begin{enumerate}[label=(\arabic*)]
\item $I_k (D) + \jepPubScript{O}_X(-Z) \subseteq I_{k-1} (D) + \jepPubScript{O}_X(-Z)$ for all $k$.
\item If $\alpha \in (0 , 1]$, then $I_k (D) = \jepPubScript{O}_X \iff k \leqslant \widetilde{\alpha}_Z - \alpha$, where 
$\widetilde{\alpha}_Z$  is the negative of the largest root of the reduced Bernstein-Sato polynomial of $Z$.
\item If $I_{k-1}(D) = \jepPubScript{O}_X$ (we say that $(X, D)$ is $(k-1)$-log canonical),
then $I_{k+1} (D) \subseteq I_k (D)$.
\item Fixing $k$, there exists a finite set of rational numbers $0 = c_0 < c_1 < \cdots < c_s < c_{s+1} = 1$ such that 
for each $0 \leqslant i \leqslant s$ and each $\alpha \in (c_i, c_{i+1}]$ we have 
$$I_k (\alpha Z)\cdot\jepPubScript{O}_Z = I_k (c_{i+1} Z)\cdot\jepPubScript{O}_Z = \text{constant}$$
and such that 
$$I_k (c_{i+1} Z)\cdot\jepPubScript{O}_Z\subsetneq I_k (c_i Z)\cdot\jepPubScript{O}_Z.$$
\end{enumerate}
\end{intro-theorem}

Going back to the description of Hodge ideals by means of log resolutions, the strictness of the Hodge filtration 
for the push-forwards of (summands of) mixed Hodge modules leads to the following local Nakano-type vanishing 
result for $\jepPubBold{Q}$-divisors:

\begin{intro-corollary}\label{intro_local_vanishing}
Let $D$ be an effective $\jepPubBold{Q}$-divisor on a smooth variety $X$ of dimension $n$, and let $f\colon Y\to X$ be a log resolution of $(X,D)$ that is an isomorphism over
$X\smallsetminus {\rm Supp}(D)$. If $E= (f^*D)_{{\rm red}}$, then
$$R^q f_*\big(\Omega_Y^p(\log E)\otimes_{\jepPubScript{O}_Y} \jepPubScript{O}_Y(-\lceil f^*D\rceil)\big)=0\quad\text{for}\quad p+q>n.$$
\end{intro-corollary}

Note that for $p =n$ this is the local vanishing for multiplier ideals \cite[Th. 9.4.1]{Lazarsfeld}, since $E - \lceil f^*D\rceil = - [(1- \varepsilon) f^*D]$ for $0 < \varepsilon \ll 1$. In general, the statement extends the case of reduced $D$ in \cite[Cor. 3]{Saito-LOG} (cf. also \cite[\S{A.5}]{Saito-MLCT}). Unlike \cite[Th. 32.1]{MP1} regarding that case, at the moment we are unable to prove this corollary via more elementary methods.

A different series of applications, given in \cite{MP3}, uses the results proved in this paper together with the relationship between Hodge ideals of $\jepPubBold{Q}$-divisors and the $V$-filtration, in order to describe the behavior of the invariant $\widetilde{\alpha}_Z$ described in Theorem \ref{other_paper} (called the \emph{minimal exponent} of $Z$). For instance, the triviality criterion proved here as Proposition 11.2 
leads to a lower bound \cite[Cor. ~D]{MP3} for  $\widetilde{\alpha}_Z$ in 
terms of invariants on a log resolution, addressing a question of Lichtin and Koll\'{a}r. Moreover, the results
in Theorem \ref{main_list} (vi) and (vii), and Corollary 11.11, lead to effective bounds and to restriction and semicontinuity statements for $\widetilde{\alpha}_Z$, in analogy with well-known properties of log canonical thresholds; for details see \cite[\S6]{MP3}.




\section*{B. Hodge ideals via log resolutions, and first properties}


Let $X$ be a smooth complex algebraic variety.
Given an effective $\jepPubBold{Q}$-divisor $D$ on $X$, our goal is to attach to $D$ ideal sheaves $I_k(D)$ for $k \geqslant 0$; when $D$ is a reduced divisor, these will coincide with the Hodge ideals in \cite{MP1}.

\section{A brief review of Hodge modules}\label{Hodge_modules}
A key ingredient for the definition of our invariants is Saito's theory of mixed Hodge modules. In what follows,
we give a brief presentation of the relevant objects, and recall a few facts that we will need.
For details, we refer to \cite{Saito-MHM}.

Given a smooth $n$-dimensional complex algebraic variety $X$, we denote by $\jepPubScript{D}_X$ the sheaf of differential operators on $X$. 
This carries the increasing filtration $F_{\jepPubBullet}\jepPubScript{D}_X$ by order of differential operators. A \emph{left} or \emph{right $\jepPubScript{D}$-module}
is a left, respectively right, $\jepPubScript{D}_X$-module, which is quasi-coherent as an $\jepPubScript{O}_X$-module. There is an equivalence 
between the categories of left and right $\jepPubScript{D}$-modules, which at the level of $\jepPubScript{O}_X$-modules is given by
$$\jepPubScript{M}\longrightarrow \jepPubScript{N}:=\jepPubScript{M}\otimes_{\jepPubScript{O}_X}\omega_X\quad\text{and}\quad \jepPubScript{N}\longrightarrow \jepPubScript{H}om_{\jepPubScript{O}_X}(\omega_X,\jepPubScript{N}).$$
For example, this equivalence maps the left $\jepPubScript{D}$-module $\jepPubScript{O}_X$ to the right $\jepPubScript{D}$-module $\omega_X$.
For a thorough introduction to the theory of $\jepPubScript{D}$-modules, we refer to \cite{HTT}.

A \emph{filtered left} (or \emph{right}) \emph{$\jepPubScript{D}$-module} 
is a $\jepPubScript{D}$-module $\jepPubScript{M}$, together with an increasing filtration $F=F_{\jepPubBullet}\jepPubScript{M}$ that is compatible with the
order filtration on $\jepPubScript{D}_X$ and which is \emph{good}, in a sense to be defined momentarily. A morphism of filtered $\jepPubScript{D}$-modules is required to be compatible with the filtrations. The equivalence between left
and right $\jepPubScript{D}$-modules extends to the categories of filtered modules, with the convention that 
$$F_{p-n}(\jepPubScript{M}\otimes_{\jepPubScript{O}_X}\omega_X)=F_p\jepPubScript{M}\otimes_{\jepPubScript{O}_X}\omega_X.$$
A filtration $F_{\jepPubBullet}\jepPubScript{M}$ on a coherent $\jepPubScript{D}$-module $\jepPubScript{M}$ is \emph{good} if the corresponding graded object 
${\rm gr}_{\jepPubBullet}^F\jepPubScript{M}:=\mathop{\textstyle\bigoplus}\limits_kF_k\jepPubScript{M}/F_{k-1}\jepPubScript{M}$ is locally finitely generated over ${\rm gr}_{\jepPubBullet}^F\jepPubScript{D}_X$. 
We note that
every coherent $\jepPubScript{D}$-module admits a good filtration, but this is far from being unique. 


We now come to the key objects in Saito's theory, the \emph{mixed Hodge modules} from \cite{Saito-MHM}. Such an object is given by the data 
$M=(\jepPubScript{M},F,\jepPubScript{P} ,\varphi, W)$, where:
\begin{enumerate}
\item[(i)] $(\jepPubScript{M},F)$ is a filtered $\jepPubScript{D}$-module, with $\jepPubScript{M}$ a holonomic left (or right) $\jepPubScript{D}$-module, with regular singularities; 
$F$ is the \emph{Hodge filtration} of $\jepPubScript{M}$.
\item[(ii)] $\jepPubScript{P}$ is a perverse sheaf of $\jepPubBold{Q}$-vector spaces on $X$.
\item[(iii)] $\varphi$ is an isomorphism between $\jepPubScript{P}_{\jepPubBold{C}} = \jepPubScript{P} \otimes_{\jepPubBold{Q}} \jepPubBold{C}$ and ${\rm DR}(\jepPubScript{M})$, i.e. 
the perverse sheaf corresponding to $\jepPubScript{M}$ via the Riemann-Hilbert correspondence.
\item[(iv)] $W$ is a finite, increasing filtration on $(\jepPubScript{M},F, \jepPubScript{P},\varphi)$, the \emph{weight filtration} of the mixed Hodge module.
\end{enumerate}
For a such an object to be a mixed Hodge module, it has to satisfy 
a complicated set of conditions of an inductive nature, which we do not discuss here.
The main reference for the basic definitions and results of this theory is \cite{Saito-MHM}; see also 
\cite{Saito-YPG} for an introduction. 


Given a mixed Hodge module $(\jepPubScript{M},F,\jepPubScript{P},\varphi, W)$, we say that the filtered $\jepPubScript{D}$-module $(\jepPubScript{M},F)$ is a \emph{Hodge $\jepPubScript{D}$-module}
(or that it
\emph{underlies} a mixed Hodge module). In fact, this is the only piece of information that we will be concerned with in this article. 
The basic example of a mixed Hodge module is ${\mathbf Q}_X^H[n]$, the trivial one. In this case, the filtered $\jepPubScript{D}$-module is the structure sheaf $\jepPubScript{O}_X$, 
with the filtration such that ${\rm gr}_p^F\jepPubScript{O}_X=0$ for all $p\neq 0$. The corresponding perverse sheaf is $\jepPubBold{Q}_X [n]$ and the weight filtration is such that 
${\rm gr}_p^W\jepPubScript{O}_X=0$ for $p\neq n$. 

The mixed Hodge modules on $X$ form an Abelian category, denoted ${\rm MHM}(X)$. Morphisms in this category are strict with respect to both the Hodge
and the weight filtration. The corresponding bounded derived category is denoted ${\bf D}^b\big({\rm MHM}(X)\big)$. 

Mixed Hodge modules satisfy Grothendieck's 6 operations formalism. The relevant fact for us is that to every morphism $f\colon X\to Y$ of 
smooth complex algebraic varieties we have a corresponding push-forward functor $f_+\colon {\bf D}^b\big({\rm MHM}(X)\big)\rightarrow {\bf D}^b\big({\rm MHM}(Y)\big)$
(this is denoted by $f_*$ in \cite{Saito-MHM}). Moreover, if
$g\colon Y\to Z$ is another such morphism, we have a functorial isomorphism $(g\circ f)_+\simeq g_+\circ f_+$. 

Regarding the push-forward functor for mixed Hodge modules, we note that on the level of $\jepPubScript{D}$-modules, it coincides with the usual $\jepPubScript{D}$-module push-forward. 
Moreover, if $f\colon X\to Y$ is proper and if we denote by ${\rm FM}(\jepPubScript{D}_X)$ the category of filtered $\jepPubScript{D}$-modules on $X$ (here it is convenient to work with right $\jepPubScript{D}$-modules), then
Saito defined in \cite{Saito-MHP} a functor
$$f_+ \colon {\bf D}^b \big({\rm FM}(\jepPubScript{D}_X)\big) \longrightarrow {\bf D}^b \big({\rm FM}(\jepPubScript{D}_Y)\big).$$
This is compatible with the usual direct image functor for right $\jepPubScript{D}$-modules and it is used to define
the push-forward between the derived categories of mixed Hodge modules at the level of filtered complexes. 
With a slight abuse of notation, if $(\jepPubScript{M},F)$ underlies a mixed Hodge module $M$ on $X$ and if $f\colon X\to Y$
is an arbitrary morphism, then we write $f_+(\jepPubScript{M},F)$ for the object in ${\bf D}^b\big({\rm FM}(\jepPubScript{D}_Y)\big)$ underlying $f_+M$.


An important feature of the push-forward of Hodge $\jepPubScript{D}$-modules with respect to proper morphisms is \emph{strictness}. 
This says that if $f\colon X\to Y$ is proper and $(\jepPubScript{M},F)$ underlies a mixed Hodge module on $X$, then 
$f_+ (\jepPubScript{M}, F)$ is strict 
as an object in ${\bf D}^b \big({\rm FM}(\jepPubScript{D}_Y)\big)$ (and moreover, each $H^i f_+ (\jepPubScript{M}, F)$ underlies a Hodge $\jepPubScript{D}_Y$-module). 
This means that the natural mapping 
\begin{equation}\label{strictness_formula}
R^i f_* \big(F_k (\jepPubScript{M} \overset{\jepPubBold{L}}{\otimes}_{\jepPubScript{D}_X} \jepPubScript{D}_{X\to Y}) \big) \longrightarrow 
R^i f_* (\jepPubScript{M} \overset{\jepPubBold{L}}{\otimes}_{\jepPubScript{D}_X} \jepPubScript{D}_{X\to Y})
\end{equation}
is injective for every $i, k\in\jepPubBold{Z}$. 
Taking $F_kH^if_+(\jepPubScript{M},F)$ to be the image of this map, we get the filtration on $H^if_+(\jepPubScript{M},F)$.

The push-forward with respect to open embeddings is more subtle. For example, suppose that $Z$ is an effective divisor
on the smooth variety $X$ and $j\colon U=X\smallsetminus Z\hookrightarrow X$ is the corresponding open immersion. 
Recall that $\jepPubScript{O}_X(*Z)$
is the push-forward $j_*\jepPubScript{O}_U$; on a suitable affine open neighborhood $V$ of a given point in $X$, this is given by localizing $\jepPubScript{O}_X(V)$ at 
an equation defining $Z\cap V$ in $V$.
$\jepPubScript{O}_X(*Z)$ has a natural left $\jepPubScript{D}$-module structure induced by the canonical $\jepPubScript{D}$-module structure on $\jepPubScript{O}_X$. In fact, as such
we have $\jepPubScript{O}_X(*Z)\simeq j_+\jepPubScript{O}_U$  (in general, for a $\jepPubScript{D}_U$-module $\jepPubScript{M}$, the $\jepPubScript{D}$-module push-forward
 $j_+\jepPubScript{M}$ agrees with $j_*\jepPubScript{M}$, with the induced $\jepPubScript{D}_X$-module structure).
We thus see that $\jepPubScript{O}_X(*Z)$ carries a canonical filtration such that the corresponding
filtered $\jepPubScript{D}$-module underlies $j_+\jepPubBold{Q}_U^H[n]$. This filtration is the one that leads to the Hodge ideals studied in \cite{MP1}.



\section{Filtered \texorpdfstring{\mbox{\normalsize$\jepPubScript{D}$}}{D}-modules associated to \texorpdfstring{\mbox{\normalsize$\jepPubBold{Q}$}}{Q}-divisors}\label{coverings}
Let $X$ be a smooth complex algebraic variety, with $\dim(X)=n$. The ideals we associate to effective $\jepPubBold{Q}$-divisors on $X$ arise from certain Hodge $\jepPubScript{D}$-modules. 
The $\jepPubScript{D}$-modules themselves have been extensively studied: these are the 
$\jepPubScript{D}$-modules attached to rational powers of functions on $X$. We proceed to recall their definition.

Consider a nonzero $h\in\jepPubScript{O}_X(X)$ and $\beta\in\jepPubBold{Q}$. We denote by $Z$ the reduced divisor on $X$ with the same support as
$H={\rm div}(h)$ and let $j\colon U=X\smallsetminus {\rm Supp}(Z)\hookrightarrow X$ be the inclusion map. 
 We consider the left $\jepPubScript{D}_X$-module 
$\jepPubScript{M}(h^{\beta})$, which is a rank 1 free $\jepPubScript{O}_X(*Z)$-module 
with generator the symbol $h^{\beta}$, on which a derivation $D$ of $\jepPubScript{O}_X$ acts by
$$D(wh^{\beta}) :=\Big(D(w) +  w\frac{\beta\cdot D(h)}{h}\Big)h^{\beta}.$$
We will denote the corresponding right $\jepPubScript{D}_X$-module by $\jepPubScript{M}_r(h^{\beta})$.
This can be described as $h^{\beta}\omega_X(*Z)$, an $\jepPubScript{O}_X$-module isomorphic to $\omega_X(*Z)$,
and such that if $x_1,\ldots,x_n$ are local coordinates, then 
$$(h^{\beta}w dx_1\cdots dx_n)\partial_i=- h^{\beta}\Big(\frac{\partial w}{\partial x_i}+ w\frac{\beta}{h}\cdot \frac{\partial h}{\partial x_i}\Big)dx_1\cdots dx_n$$
for every $i$ with $1\leqslant i\leqslant n$. 



\begin{remark}\label{case_ell_1}
When $\beta\in\jepPubBold{Z}$, we have a canonical isomorphism of left $\jepPubScript{D}_X$-modules
\begin{equation}\label{eq_case_ell_1}
\jepPubScript{M}(h^{\beta})\simeq \jepPubScript{O}_X(*Z),\,\,\,\,\,\,wh^{\beta}\longmapsto w \cdot h^{\beta},
\end{equation}
where on the localization $\jepPubScript{O}_X(*Z)$ we consider the natural $\jepPubScript{D}_X$-module structure induced from $\jepPubScript{O}_X$. Note that $\jepPubScript{O}_X(*Z)$
is also the $\jepPubScript{D}$-module push-forward $j_+ \jepPubScript{O}_U$.
\end{remark}

\begin{remark}\label{eq_renormalization_0}
For every positive integer $m$, we have a canonical isomorphism of  left $\jepPubScript{D}_X$-modules
$$\jepPubScript{M}(h^{\beta})\simeq\jepPubScript{M}\big((h^m)^{\beta /m}\big),\,\,\,\,\,\,wh^{\beta}\longmapsto w(h^m)^{\beta/m}.$$
\end{remark}


\begin{remark}\label{twist_individual}
We can define, more generally, left $\jepPubScript{D}$-modules $\jepPubScript{M}(h_1^{\beta_1}\cdots h_r^{\beta_r})$, for 
nonzero regular functions $h_1,\ldots,h_r\in\jepPubScript{O}_X(X)$ and rational numbers $\beta_1,\ldots,\beta_r$. 
If $\ell_i$ are positive integers such that
$\beta_i/\ell_i=\beta$ for all $i$ and if $h=\prod_ih_i^{\ell_i}$, then we have an isomorphism of left $\jepPubScript{D}_X$-modules
$$\jepPubScript{M}(h_1^{\beta_1}\cdots h_r^{\beta_r})\simeq \jepPubScript{M}(h^{\beta}).$$
\end{remark}

\begin{remark}\label{twist_positive}
If $r$ is an integer, then we have an isomorphism of left $\jepPubScript{D}_X$-modules
$$\jepPubScript{M}(h^{\beta})\longrightarrow \jepPubScript{M}(h^{r + \beta}), \,\,\,\,\,\,wh^{\beta}\longmapsto (wh^{-r}) h^{r+\beta}.$$
\end{remark}

\medskip

Let now $D$ be an effective $\jepPubBold{Q}$-divisor on $X$. We denote by $Z$ the reduced divisor with the same support as $D$. 
As above, we put $U=X\smallsetminus Z$ and let $j\colon U\hookrightarrow X$ be the inclusion map. 
We first assume that we can write $D=\alpha\cdot {\rm div}(h)$ for some nonzero $h\in\jepPubScript{O}_X(X)$ and $\alpha\in\jepPubBold{Q}_{>0}$ (this is of course 
always the case locally). To this data we can associate the $\jepPubScript{D}_X$-module $\jepPubScript{M}(h^{-\alpha})$; later it will be more convenient to 
consider equivalently (according to Remark \ref{twist_positive}) the $\jepPubScript{D}_X$-module $\jepPubScript{M}(h^{1-\alpha})$. 
This depends on the choice of $h$; however, if we replace $h$ by $h^m$ and $\alpha$ by $\alpha/m$, for some positive integer $m$,
the $\jepPubScript{D}$-module does not change (see Remark~\ref{eq_renormalization_0}). In particular, we may always assume that $\alpha=1/\ell$,
for a positive integer $\ell$.




\begin{remark}\label{rmk_add_eff_div}
Suppose that $D'$ is a $\jepPubBold{Q}$-divisor with the same support as $D$ and such that $D-D'={\rm div}(u)$, for some $u\in\jepPubScript{O}_X(X)$.
Suppose that we can write
$D'=(1/\ell)\cdot {\rm div}(h')$ for some $h'\in\jepPubScript{O}_X(X)$ and some positive integer $\ell$.
In this case we can also write
$D=(1/\ell)\cdot {\rm div}(h)$, where $h=u^{\ell}{h'}$, and  we
have an isomorphism of  $\jepPubScript{D}_X$-modules
\begin{equation}\label{isom_add_1}
\jepPubScript{M}(h^{-1/\ell})\longrightarrow \jepPubScript{M}({h'}^{-1/\ell}), \,\,\,\,\,\,gh^{-1/\ell} \longmapsto gu^{-1}{h'}^{-1/\ell}.
\end{equation}
\end{remark}

\medskip

Our first goal is to show that $\jepPubScript{M}(h^{-\alpha})$ is canonically a filtered $\jepPubScript{D}_X$-module. 
Let $\ell$ be a positive integer such that $\ell\alpha\in\jepPubBold{Z}$.
Consider the finite \'{e}tale map $p\colon V\to U$, where $V={\bf Spec}\,\jepPubScript{O}_U[y]/(y^{\ell}-h^{-\ell\alpha})$.
Note that this fits in a Cartesian diagram
\begin{equation}\label{cart_diag1}
{\def\labelstyle{\textstyle}\xymatrix{V\ar[r]\ar[d]_{p}&W\ar[d]^{q}\\U\ar[r]_{j}&X,}}
\end{equation}
in which
$$W={\bf Spec}\,\jepPubScript{O}_X[z]/(z^\ell-h^{\ell\alpha}),$$
such that the map $V\to W$ pulls $z$ back to $y^{-1}=y^{\ell-1}h^{\ell\alpha}$. 


\begin{lemma}\label{lem_decomp1}
We have an isomorphism of left $\jepPubScript{D}_X$-modules
\begin{equation}\label{eq1_lem_decomp1}
j_+p_+\jepPubScript{O}_V\simeq\mathop{\textstyle\bigoplus}\limits_{i=0}^{\ell-1}\jepPubScript{M}(h^{-i\alpha}),
\end{equation}
with the convention that the first summand is $\jepPubScript{O}_X(*Z)$.
\end{lemma}

\begin{proof}
Since $p$ is finite \'{e}tale, it follows that we have a canonical isomorphism $\tau\colon p^*\jepPubScript{D}_U\simeq\jepPubScript{D}_V$, 
and for every $\jepPubScript{D}_V$-module $\jepPubScript{M}$ we have $p_+\jepPubScript{M}\simeq p_*\jepPubScript{M}$,
with the action of $\jepPubScript{D}_U$ induced via the isomorphism $\tau$.

By mapping $gy^i$ to $gh^{-i\alpha}$, where $g$ is a section of $\jepPubScript{O}_X$ and $0\leqslant i\leqslant\ell-1$, we obtain an isomorphism 
of $\jepPubScript{O}_X$-modules as in
(\ref{eq1_lem_decomp1}). In order to see that this is an isomorphism of $\jepPubScript{D}_X$-modules, consider a local derivation
$D$ of $\jepPubScript{O}_X$ and note that since $y^{\ell}=h^{-\ell\alpha}$, by identifying $D$ with its pull-back to $V$ we have
$$D (y^i)=iy^{i-1}D(y)=-i\alpha y^i\frac{D(h)}{h},$$
which via our map corresponds to $D (h^{-i\alpha})$.
This implies the assertion.
\end{proof}

It follows from the lemma that the right-hand side of (\ref{eq1_lem_decomp1}) is the $\jepPubScript{D}$-module corresponding to the 
mixed Hodge module push-forward $(j\circ p)_+\jepPubBold{Q}_V^H[n]$. In particular, it 
 carries a canonical structure of filtered $\jepPubScript{D}$-module. 


\begin{remark}\label{rem_replace_by_multiple}
Let us see what happens if we replace $\ell$ by a multiple $m\ell$. Let $p_{\ell}\colon V_{\ell}\to U$ and $p_{m\ell}\colon V_{m\ell}\to U$ be the corresponding \'{e}tale covers.
Note that $$V_{m\ell}={\bf Spec}\,\jepPubScript{O}_U[y]/(y^{\ell m}-h^{-\ell m\alpha})$$ decomposes as a disjoint union of $m$ copies of $V_{\ell}$, and thus we have an isomorphism of filtered $\jepPubScript{D}_X$-modules (and a corresponding isomorphism of mixed Hodge modules)
\begin{equation}\label{eq_rem_replace_by_multiple}
j_+(p_{m\ell})_+\jepPubScript{O}_{V_{m\ell}}\simeq 
\big(j_+(p_{\ell})_+\jepPubScript{O}_{V_{\ell}}\big)^{\oplus m}.
\end{equation}
If $\eta$ is a primitive root of $1$ of order $\ell m$, and if on each side of 
(\ref{eq_rem_replace_by_multiple}) we consider the decompositions 
(\ref{eq1_lem_decomp1}), then the isomorphism maps 
$$h^{-i\alpha} \longmapsto \big(\eta^{is}h^{-c\ell\alpha}\cdot h^{-d\alpha}\big)_{0\leqslant s\leqslant m-1},$$
where we write $i=\ell c+d$, with $0\leqslant c\leqslant m-1$ and $0\leqslant d\leqslant\ell-1$.
\end{remark}

We can interpret the isomorphism in (\ref{eq1_lem_decomp1}) in terms of a suitable $\mu_\ell$-action, where $\mu_{\ell}$
is the group of $\ell$-th roots of $1$ in $\jepPubBold{C}^*$. Note that we have a natural action of $\mu_{\ell}$ on $W$ such that
via the corresponding action on $\jepPubScript{O}_W$, an element $\lambda\in\mu_{\ell}$ maps $z^i$ to $\lambda^iz^i$. 
If we let $\mu_{\ell}$ act trivially on $X$, then $q$ is an equivariant morphism (in fact, $q$ is the quotient morphism with respect
to the $\mu_{\ell}$-action). It is clear that $q^{-1}(Z)$ is fixed by the $\mu_{\ell}$-action and we have an induced $\mu_{\ell}$-action
on $W\smallsetminus q^{-1}(Z)=V$. This in turn induces a $\mu_{\ell}$-action on $j_+p_+\jepPubScript{O}_V$ and the isomorphism in
(\ref{eq1_lem_decomp1}) corresponds to the isotypic decomposition of $j_+p_+\jepPubScript{O}_V$, such that every 
$\lambda\in\mu_{\ell}$ acts on $\jepPubScript{M}(h^{-i\alpha})$ by multiplication with $\lambda^{-i}$. 


\begin{lemma}\label{decomp_filtration}
The filtration on $j_+p_+\jepPubScript{O}_V$ is preserved by the $\mu_{\ell}$-action. Therefore we have an induced filtration on each 
$\jepPubScript{M}(h^{-i\alpha})$ such that (\ref{eq1_lem_decomp1}) is an isomorphism of filtered $\jepPubScript{D}$-modules.
\end{lemma}


\begin{proof}
One way to see this is by using a suitable equivariant resolution of $W$. Let $W'$ be the disjoint union of the irreducible components of $W$ and
$q'\colon W'\to W$ the canonical morphism. It is clear that the $\mu_{\ell}$-action on $W$ induces an action on $W'$ such that $q'$ is equivariant.
Since $V$ is contained in the smooth locus of $W$, it has an open immersion into $W'$. 
We use equivariant resolution of singularities to construct a $\mu_{\ell}$-equivariant morphism 
$\varphi\colon Y\to W'$ that is an isomorphism over $V$ and such that $(q\circ q'\circ\varphi)^*(Z)$ is a divisor with simple normal crossings.  
Let $g=q\circ q'\circ \varphi$. If $E$ is the reduced, effective divisor supported on $g^{-1}(Z)$, then 
we have an isomorphism of filtered $\jepPubScript{D}$-modules (induced by a corresponding isomorphism of mixed Hodge modules)
\begin{equation}\label{eq_ism_filtered}
j_+p_+\jepPubScript{O}_V\simeq g_+\widetilde{j}_+\jepPubScript{O}_V\simeq g_+\jepPubScript{O}_Y(*E),
\end{equation}
where $\widetilde{j}\colon Y\smallsetminus {\rm Supp}(E)\hookrightarrow Y$ is the inclusion map.

We can deduce the assertion in the lemma from an explicit
computation of the filtration on $j_+p_+\jepPubScript{O}_V$ via the isomorphism (\ref{eq_ism_filtered}), as follows. First, since we deal with
$\jepPubScript{D}$-module push-forward, it is more convenient to work with right $\jepPubScript{D}$-modules. We will thus compute $g_+\omega_Y(*E)$, where
$\omega_Y(*E)$ is the filtered right $\jepPubScript{D}$-module corresponding to $\jepPubScript{O}_Y(*E)$. 

Since $E$ is a simple normal crossing divisor, $\omega_Y(*E)$ has a resolution by a complex $C^\jepPubBullet$ 
of filtered right $\jepPubScript{D}_Y$-modules
$$0\longrightarrow C^{-n}\longrightarrow\cdots\longrightarrow C^0\longrightarrow 0,$$
where $C^{i}=\Omega_Y^{i+n}(\log E)\otimes_{\jepPubScript{O}_Y}\jepPubScript{D}_Y$, with the filtration given by
$$F_{k-n}C^i=\Omega_Y^{i+n}(\log E)\otimes_{\jepPubScript{O}_Y}F_{k+i}\jepPubScript{D}_Y.$$ 
For a description of the maps in this complex, see the beginning of Section~\ref{complex_for_SNC} below;
a proof
of the fact that it resolves $\omega_Y(*E)$ is given in \cite[Prop.~3.1]{MP1}. We can thus compute  $F_kg_+\omega_Y(*E)$ as the image of the injective map
$$R^0g_*\big(F_k(C^{\jepPubBullet}\otimes_{\jepPubScript{D}_Y}\jepPubScript{D}_{Y\to X})\big)\longrightarrow R^0g_*(C^{\jepPubBullet}\otimes_{\jepPubScript{D}_Y}\jepPubScript{D}_{Y\to X})=g_+\omega_Y(*E).$$
Since $g$ is equivariant and the action of $\mu_{\ell}$ on $Y$ induces an action on $E$ (in fact, it fixes $E$), the above description implies that each
$F_kg_+\omega_Y(*E)$ is preserved by the $\mu_{\ell}$-action.
\end{proof}

\begin{remark}\label{rmk_filtration_first_piece}
We note that the filtration on $j_+p_+\jepPubScript{O}_V$ induces the canonical filtration on the first summand $\jepPubScript{O}_X(*Z)$.
Indeed, on $U$ we have a morphism of mixed Hodge modules $\jepPubBold{Q}_U^H[n]\to p_+\jepPubBold{Q}_V^H[n]$. Applying $j_+$ and only considering the underlying filtered $\jepPubScript{D}$-modules, we obtain a morphism $j_+\jepPubScript{O}_U\to j_+p_+\jepPubScript{O}_V$, which is an isomorphism onto the first summand.
\end{remark}

\begin{definition}
Given $\alpha>0$, choose $\ell\geqslant 2$ such that $\ell\alpha\in\jepPubBold{Z}$. In this case $\jepPubScript{M}(h^{-\alpha})$ appears as the second summand in the decomposition
(\ref{eq1_lem_decomp1}). We define the filtration 
$$F_k \jepPubScript{M}(h^{-\alpha}) \,\,\,\,\,\,\text{for} \,\,\,\,\,\, k \geqslant 0$$ 
to be the filtration induced from the canonical filtration on $j_+p_+\jepPubScript{O}_V$. 
It is straightforward to see, using the discussion in Remark~\ref{rem_replace_by_multiple} that this filtration does not change if we replace $\ell$
by a multiple; therefore it is independent of $\ell$. Moreover, we note that if $\alpha$ is an integer, using the same Remark~\ref{rem_replace_by_multiple}, 
the isomorphism $\jepPubScript{M}(h^{-\alpha})\simeq\jepPubScript{O}_X(*Z)$ is an isomorphism of filtered $\jepPubScript{D}$-modules.
\end{definition}

In this definition, a priori different covers have to be considered for each of the summands $\jepPubScript{M}(h^{-i\alpha})$. However, 
we have:

\begin{lemma}\label{lem_decomp1_v2}
With the filtration defined above, the isomorphism (\ref{eq1_lem_decomp1}) is an isomorphism of filtered $\jepPubScript{D}$-modules. 
\end{lemma}
\begin{proof}
By Lemma~\ref{decomp_filtration}, we only need to show that for every $i$ with $0\leqslant i\leqslant \ell-1$, the filtration induced on $\jepPubScript{M}(h^{-i\alpha})$ by that on $j_+p_+\jepPubScript{O}_V$ coincides with the 
one given in the above definition. For $i=0$ this follows from Remark~\ref{rmk_filtration_first_piece}. If $i>0$, consider the cover used to define the filtration on 
$\jepPubScript{M}(h^{-i\alpha})$, namely 
$$p'\colon V'={\bf Spec}\jepPubScript{O}_U[y]/(y^{\ell}-h^{i\alpha\ell})\longrightarrow U.$$ 
Note that we have a finite morphism
$\psi\colon V\to V'$ of varieties over $U$, that pulls-back $y$ to $y^i$. We have a canonical morphism of mixed Hodge modules
$\jepPubBold{Q}_{V'}^H[n]\to \psi_+\jepPubBold{Q}_V^H[n]$. Applying $j_+p'_+$ and passing to the underlying filtered $\jepPubScript{D}$-modules, we obtain a morphism of filtered 
$\jepPubScript{D}$-modules $j_+p'_+\jepPubScript{O}_{V'}\to j_+p_+\jepPubScript{O}_V$ that is the identity on the summand $\jepPubScript{M}(h^{-i\alpha})$. This proves our claim.
\end{proof}

\begin{remark}\label{same_filtration_rescaling}
It is clear from definition that for every $\alpha>0$ and every positive integer $m$, the isomorphism
$$\jepPubScript{M}(h^{-\alpha})\longrightarrow \jepPubScript{M}\big((h^m)^{-\alpha/m}\big), \quad gh^m\longmapsto g (h^m)^{-\alpha/m}$$
is an isomorphism of filtered $\jepPubScript{D}$-modules. 
\end{remark}


\begin{remark}\label{rem_isom_comp_with_filtrations}
In the setting of Remark~\ref{rmk_add_eff_div},
the isomorphism (\ref{isom_add_1}) is an isomorphism of filtered $\jepPubScript{D}_X$-modules.
This is clear if $\ell=1$, hence we assume $\ell\geqslant 2$.
Let
$p\colon V\to U$ and $p'\colon V'\to U$ be the canonical projections, where
$$V={\bf Spec}\,\jepPubScript{O}_U[y]/(y^{\ell}-h)\quad\text{and}
\quad V'={\bf Spec}\,\jepPubScript{O}_U[y]/(y^{\ell}-h').$$
We have an isomorphism $\varphi\colon V'\to V$ of schemes over $U$, where $\varphi^*(y)=uy$. 
This induces an isomorphism of filtered $\jepPubScript{D}_X$-modules
$$j_+p_+\jepPubScript{O}_V\simeq j_+p'_+\jepPubScript{O}_{V'},$$
which via the identifications given by Lemma~\ref{lem_decomp1} is the direct sum
$$\mathop{\textstyle\bigoplus}\limits_{i=0}^{\ell-1}\jepPubScript{M}(h^{-i/\ell})\simeq 
\mathop{\textstyle\bigoplus}\limits_{i=0}^{\ell-1}\jepPubScript{M}({h'}^{-i/\ell})$$
of the isomorphisms (\ref{isom_add_1}). For $i=1$, we obtain our assertion.
\end{remark}

A special case of the above remark implies that for every $\alpha>0$ the isomorphism
$$\jepPubScript{M}(h^{-\alpha})\longrightarrow\jepPubScript{M}(h^{-\alpha-1}),\quad gh^{-\alpha}\longmapsto (gh)h^{-\alpha-1}$$
is an isomorphism of filtered $\jepPubScript{D}$-modules. We use this to put a structure of filtered
$\jepPubScript{D}$-module on $\jepPubScript{M}(h^{\beta})$ for every $\beta\in\jepPubBold{Q}$, such that for every $r\in\jepPubBold{Z}$,
we have an isomorphism of filtered $\jepPubScript{D}$-modules
$$\jepPubScript{M}(h^{\beta})\longrightarrow\jepPubScript{M}(h^{\beta-r}),\quad gh^{\beta}\longmapsto (gh^r)h^{\beta-r}.$$
For example, we have have an isomorphism of filtered $\jepPubScript{D}$-modules $\jepPubScript{M}(h^0)\simeq \jepPubScript{O}_X(*Z)$.


\begin{remark}\label{rem_indep_equations1}
Suppose that $h, \overline{h} \in\jepPubScript{O}_X(X)$ are nonzero, and $\alpha,\overline{\alpha}\in\jepPubBold{Q}_{>0}$ are 
such that we have the equality of $\jepPubBold{Q}$-divisors
$$\alpha \cdot {\rm div}(h)=\overline{\alpha}\cdot {\rm div}(\overline{h}).$$
Let $\ell$ be  a positive integer such that $\ell\alpha,\ell\overline{\alpha}\in\jepPubBold{Z}$.
In this case there is $g\in\jepPubScript{O}_X^*(X)$ such that $h^{\ell\alpha}=g{\overline{h}}^{\ell\overline{\alpha}}$.
 Suppose now that there exists $G\in\jepPubScript{O}_X(X)$ such that $G^{\ell}=g$. (For example, this holds after pulling-back to the
\'{e}tale cover $\mathbf{Spec}\,\jepPubScript{O}_X[z]/(z^{\ell}-g)$.)
In this case we have an isomorphism of filtered $\jepPubScript{D}_X$-modules
$$\Phi\colon \jepPubScript{M}(h^{-\alpha})\longrightarrow \jepPubScript{M}(\overline{h}^{-\overline{\alpha}})$$
given by
$$\Phi(wh^{-\alpha})=wG^{-1}\overline{h}^{-\overline{\alpha}}.$$
Indeed, this follows from the definition of the filtrations and the isomorphism of schemes over $U$
$$\varphi\colon {\bf Spec} \jepPubScript{O}_U[y]/(y^{\ell}-\overline{h}^{-\ell\overline{\alpha}})\longrightarrow {\bf Spec} \jepPubScript{O}_U[y]/(y^{\ell}-h^{-\ell\alpha})$$
that pulls-back $y$ to $G^{-1}y$. 
\end{remark}

\begin{remark}\label{rem_behavior_etale}
It is clear that the filtration on $\jepPubScript{M}(h^{-\alpha})$ is compatible with restriction to open subsets.
More generally, it is compatible with smooth pullback, as follows.
Suppose that $h\in\jepPubScript{O}_X(X)$ is nonzero and $\alpha \in \jepPubBold{Q}$.
If $\varphi\colon Y\to X$ is a smooth morphism and $g=h\circ \varphi$, then there is an isomorphism of $\jepPubScript{D}_Y$-modules
$$\jepPubScript{M}(g^{-\alpha})\simeq\varphi^*\jepPubScript{M}(h^{-\alpha}),$$
such that for every $k$ we have
$$F_k\jepPubScript{M}(g^{-\alpha})\simeq\varphi^*F_k\jepPubScript{M}(h^{-\alpha}).$$
Indeed, choose $\ell\geqslant 2$ such that $\ell\alpha\in\jepPubBold{Z}$ and consider the 
Cartesian diagram
$$
{\def\labelstyle{\textstyle}\xymatrix{V_Y\ar[r]^{p_Y}\ar[d]_{\psi}&U_Y\ar[r]^{j_Y}\ar[d]&Y\ar[d]^{\textstyle\varphi}\\V\ar[r]_{p}&U\ar[r]_{j}&X,}}
$$
where $j$ and $p$ are as in Lemma~\ref{lem_decomp1} and $j_Y$ and $p_Y$ are the corresponding morphisms for $Y$ and $g$. 
Note that we have a base-change theorem that gives
\begin{equation}\label{eq_rem_behavior_etale}
\varphi^!j_+p_+\jepPubBold{Q}^H_V[n]\simeq (j_Y)_+(p_Y)_+\psi^!\jepPubBold{Q}^H_V[n]
\end{equation}
(see \cite[(4.4.3)]{Saito-MHM}).
Moreover, since $\varphi$ is smooth, if $d=\dim(Y)-\dim(X)$, then for every filtered $\jepPubScript{D}$-module 
$(\jepPubScript{M},F)$ underlying a mixed Hodge module $M$, the filtered $\jepPubScript{D}$-module
underlying $\varphi^!M$ is $(\varphi^*\jepPubScript{M},F)[d]$, where $F_k(\varphi^*\jepPubScript{M})=\varphi^*(F_{k}\jepPubScript{M})$
(see \cite[3.5]{Saito-MHP}).
This also applies to $\psi$; in particular, we have $\psi^!\jepPubBold{Q}_V^H[n]\simeq\jepPubBold{Q}_{V_Y}^H[n+2d]$.
By decomposing both sides of (\ref{eq_rem_behavior_etale}) with respect to the $\mu_{\ell}$-action, we
obtain our assertion.
\end{remark}










\section{The case of smooth divisors}\label{smooth}


Our goal now is to describe the filtrations introduced in the previous section when $Z$ is a smooth divisor. We will then use this to define
Hodge ideals for arbitrary $\jepPubBold{Q}$-divisors. The key result in the smooth case is the following:

\begin{lemma}\label{lem_decomp2}
Let 
 $$\psi\colon Y=\operatorname{Spec} \jepPubBold{C}[t]\longrightarrow X=\operatorname{Spec} \jepPubBold{C}[x]$$ be
  the map given by $\psi^*(x)=t^{\ell}$. If $Z$ is the divisor on $Y$ defined by
$t$, then
we have
an isomorphism of filtered $\jepPubScript{D}_X$-modules
$$\psi_+\jepPubScript{O}_{Y}(*Z)\simeq\mathop{\textstyle\bigoplus}\limits_{j=0}^{\ell-1}\jepPubScript{M}_j,$$
where 
$\jepPubScript{M}_j\simeq\jepPubScript{D}_X/\jepPubScript{D}_X(\partial_x x - j/\ell)$ 
and $F_k\jepPubScript{M}_j$ is generated over $\jepPubScript{O}_X$ by the classes of $1,\partial_x,\ldots,\partial_x^k$.
Moreover, if we consider on $Y$ the $\mu_{\ell}$-action such that every $\lambda\in\mu_{\ell}$ maps $t$ to $\lambda t$, then
$\jepPubScript{M}_j$ is
the component of $\psi_+\jepPubScript{O}_{Y}(*Z)$
on which every  $\lambda\in\mu_{\ell}$ acts by multiplication with $\lambda^j$.
\end{lemma}

\begin{proof}
As usual, it is easier to do the computation for the filtered right $\jepPubScript{D}$-module $\omega_Y(*Z)$ corresponding to $\jepPubScript{O}_Y(*Z)$. 
Note that this is filtered quasi-isomorphic to the complex
$$A^{\jepPubBullet}:\,\,0\longrightarrow \jepPubScript{D}_Y\xrightarrow{\;\textstyle w\;} \omega_Y(Z)\otimes_{\jepPubScript{O}_Y}\jepPubScript{D}_Y\longrightarrow 0,$$
placed in degrees $-1$ and $0$, where $w(1)=(dt/t)\otimes t\partial_t$; see e.g. \cite[Prop. 3.1]{MP1}. 
Since $\psi$ is finite, the functor $\psi_*$ is exact on quasi-coherent $\jepPubScript{O}_Y$-modules, hence
$\psi_+\omega_Y(*Z)$ is computed by the $0$-th cohomology of the complex
$$B^{\jepPubBullet}=\psi_*(A^{\jepPubBullet}\otimes_{\jepPubScript{D}_Y}\jepPubScript{D}_{Y\to X}),$$
with the obvious induced filtration. The definition of $w$
immediately implies that $w\otimes 1_{\jepPubScript{D}_{Y\to X}}$ is injective. 
Note that $dt/t=(1/\ell)dx/x$ and $t\partial_t=\ell x\partial_x$.

In order to describe the complex $B^{\jepPubBullet}$, note that any element of $B^{-1}$ can be uniquely written as $\sum_{j=0}^{\ell-1}t^jP_j$,
with $P_j\in \jepPubScript{D}_X$. Similarly, any element in $B^0$ can be uniquely written as $\sum_{j=0}^{\ell-1}t^j(dx/x)Q_j$, with $Q_j\in\jepPubScript{D}_X$.
Moreover, if $\tau$ is the differential in $B^{\jepPubBullet}$, then
$$\tau\bigg(\sum_{j=0}^{\ell-1}t^jP_j\bigg)=\sum_{j=0}^{\ell-1}t^j\frac{dx}{x}(x\partial_x+j/\ell)P_j,$$
where we use the fact that $t\partial_tt^j=t^jt\partial_t+jt^j$. In other words, we have have an eigenspace decomposition 
$$B^{\jepPubBullet}\simeq \mathop{\textstyle\bigoplus}\limits_{j=0}^{\ell-1}B^{\jepPubBullet}_j,$$
where $B^{\jepPubBullet}_j$ is identified with the complex
$$0\longrightarrow\jepPubScript{D}_X\longrightarrow \jepPubScript{D}_X\longrightarrow 0,$$
with the differential mapping $P$ to $(x\partial_x+j/\ell)P$.
It follows that $B^{\jepPubBullet}$ is filtered quasi-isomorphic to
$$\mathop{\textstyle\bigoplus}\limits_{j=0}^{\ell-1}\jepPubScript{D}_X/(x\partial_x+j/\ell)\jepPubScript{D}_X,$$
where the filtration on the $j$-th component is such that
$$F_{k-1}\left(\jepPubScript{D}_X/(x\partial_x+j/\ell)\jepPubScript{D}_X\right)$$
is the $\jepPubScript{O}_X$-submodule generated by the classes of $1,\partial_x,\ldots,\partial_x^k$.
Moreover, every $\lambda\in\mu_{\ell}$ acts on the $j^{\rm th}$ factor in the above decomposition
by multiplication with $\lambda^j$. 

The assertion in the lemma now follows immediately from the explicit description
of the equivalence between the categories of left and right $\jepPubScript{D}$-modules on $X={\mathbf A}^1$. 
Indeed, recall that if $\tau$ is the $\jepPubBold{C}$-linear endomorphism of the Weyl algebra 
$\Gamma({\mathbf A}^1,\jepPubScript{D}_{{\mathbf A}^1})$
such that $\tau(PQ)=\tau(Q)\cdot\tau(P)$ for all $P$ and $Q$, and such that $\tau(t)=t$ and $\tau(\partial_t)=-\partial_t$, then the left
$\jepPubScript{D}$-module $N$ corresponding to a right $\jepPubScript{D}$-module $M$ is isomorphic to $M$ itself, with scalar multiplication given via the map $\tau$.
Moreover, for filtered $\jepPubScript{D}$-modules, via this isomorphism $F_kN$ corresponds to $F_{k-1}M$.
In particular, we see that if $M=\jepPubScript{D}_X/P\cdot\jepPubScript{D}_X$, then $N\simeq \jepPubScript{D}_X/\jepPubScript{D}_X\cdot \tau(P)$, and we obtain the statement.
\end{proof}

In what follows, we denote by $\lceil\alpha\rceil$ the smallest integer that is $\geqslant\alpha$. For a $\jepPubBold{Q}$-divisor $D=\sum_{i=1}^ra_iD_i$,
we put $\lceil D\rceil=\sum_{i=1}^r\lceil a_i\rceil D_i$. 

\begin{corollary}\label{smooth_case}
If $h\in \jepPubScript{O}_X(X)$ is nonzero and such that the support $Z$ of ${\rm div}(h)$ is smooth (possibly disconnected), then
for every $\alpha\in\jepPubBold{Q}_{> 0}$ the filtration on 
$\jepPubScript{M}(h^{- \alpha})$ is given by
$$F_k\jepPubScript{M}(h^{-\alpha})=\jepPubScript{O}_X\big((k+1)Z-\lceil D\rceil\big)h^{-\alpha}\quad \text{if}\quad k\geqslant 0,$$
where $D=\alpha\cdot{\rm div}(h)$, 
and $F_k\jepPubScript{M}(h^{-\alpha})=0$ if $k<0$.
\end{corollary}

\begin{proof}
We first reduce to the case when $Z={\rm div}(h)$. We can check the assertion in the proposition locally, hence we may assume that
$Z={\rm div}(g)$, for some $g\in\jepPubScript{O}_X(X)$, and $h=ug^m$, for some $u\in\jepPubScript{O}_X^*(X)$. Furthermore, by Remark~\ref{rem_behavior_etale},
it is enough to prove the assertion after passing to a surjective \'{e}tale cover, hence we may assume that $u=v^m$ for some $v\in\jepPubScript{O}_X^*(X)$.
After replacing $g$ by $vg$, we may thus assume that $h=g^m$. In this case we have an isomorphism
of filtered $\jepPubScript{D}$-modules $\jepPubScript{M}(h^{-\alpha})\simeq\jepPubScript{M}(g^{-m\alpha})$, hence we may and will assume that
${\rm div}(h)=Z$.


We consider the smallest positive integer $\ell$ such that $m:=\ell\alpha\in\jepPubBold{Z}$.
If $\ell=1$, then the assertion follows from the formula for the filtration on $\jepPubScript{O}_X(*Z)$ when $Z$ is smooth; see \cite[Prop. 8.2]{MP1}. Therefore from now on we assume $\ell>1$.

The morphism $h\colon X\to\jepPubBold{A}^1$ is smooth over some open neighborhood of $0$. Using Remark~\ref{rem_behavior_etale}, we see that in order to prove the corollary,
we may assume that $X=\jepPubBold{A}^1$ and $h=x$, the standard coordinate on $\jepPubBold{A}^1$.
Consider the Cartesian diagram 
$$
{\def\labelstyle{\textstyle}\xymatrix{V\ar[r]^{j_0}\ar[d]_{p}&W\ar[d]^{g}\\U\ar[r]_{j}&X,}}
$$
where 
$$j_0\colon V=\operatorname{Spec} \jepPubBold{C}[x,x^{-1},y]/(y^{\ell}-x^{-m})\longrightarrow W=\operatorname{Spec} \jepPubBold{C}[x,z]/(z^{\ell}-x^m),\quad j_0^*(z)=y^{-1}.$$
Let
 $\varphi\colon\widetilde{W}= \operatorname{Spec} \jepPubBold{C}[t] \to W$ be the normalization, given by 
$$\varphi^*(x)=t^{\ell} \,\,\,\, \text{and} \,\,\,\, \varphi^*(z)=t^m.$$ 
(Here we use that $\ell$ and $m$ are relatively prime.) 
Note that $\varphi$ is an isomorphism over $V$, hence we have an open embedding $\iota\colon V\hookrightarrow \widetilde{W}$, with complement the smooth divisor $T$ defined by $t$
(in fact, if $a$ and $b$ are integers such that $am+b\ell=1$, then $\iota^*(t)=y^{-a}x^b$). 
We thus have 
$$j_+p_+\jepPubScript{O}_V\simeq \psi_+\iota_+\jepPubScript{O}_V\simeq\psi_+\jepPubScript{O}_{\widetilde{W}}(*T),$$
where $\psi=g\circ\varphi$. We apply Lemma~\ref{lem_decomp2} for $\psi$. 
Note that $\varphi$ is a $\mu_{\ell}$-equivariant morphism if we let each $\lambda\in\mu_{\ell}$ act on $t$ by multiplication with $\lambda^a$.
By considering the behavior with respect to the
$\mu_{\ell}$-action, we see that
 in the decomposition given by the lemma, we have
$\jepPubScript{M}_j\simeq\jepPubScript{M}(x^{-\alpha})$ if and only if $ja\equiv -1$ (mod $\ell$), that is,
$j\equiv -m$ (mod $\ell$). 

Suppose first that $\alpha<1$, in which case the condition for $j$ is that $j=\ell-m$. As a reality check, note that we indeed have an isomorphism
$$\jepPubScript{D}_X/\jepPubScript{D}_X(\partial_xx -(\ell-m)/\ell)\simeq \jepPubScript{M}(x^{-\alpha})$$
that maps the class of $1$ to $x^{-\alpha}$. 
The formula for the filtration on $\jepPubScript{M}(h^{\alpha})$ now follows from Lemma~\ref{lem_decomp2}.
When $\alpha>1$, we put $m=\lceil\alpha\rceil-1$, and use the fact from Remark \ref{twist_positive}, namely that we have an isomorphism of filtered modules
$$\jepPubScript{M}(x^{-\alpha})\longrightarrow\jepPubScript{M}(x^{-\alpha+m}),\quad gx^{-\alpha}\longmapsto (gx^{-m})x^{-\alpha+m},$$
to reduce the assertion to the case $\alpha\in (0,1)$.
This completes the proof of the corollary. 
\end{proof}




\section{Definition of Hodge ideals for \texorpdfstring{\mbox{\normalsize$\jepPubBold{Q}$}}{Q}-divisors}\label{definition}
In general, we obtain an upper bound for the terms in the filtration on $\jepPubScript{M}(h^{-\alpha})$
by restricting to the open subset where the support of ${\rm div}(h)$ is smooth, as follows.



\begin{proposition}\label{prop_for_def}
Given a nonzero $h\in\jepPubScript{O}_X(X)$ and a positive rational number $\alpha$, 
for every $k\geqslant 0$ we have 
$$F_k\jepPubScript{M}(h^{-\alpha})\subseteq \jepPubScript{O}_X\big((k+1)Z-\lceil D\rceil \big)h^{-\alpha},$$
where $D=\alpha\cdot {\rm div}(h)$ and $Z={\rm Supp}(D)$,
while $F_k\jepPubScript{M}(h^{-\alpha})=0$ for $k<0$.
\end{proposition}

\begin{proof}
Let $\iota\colon X_0\to X$ be an open immersion such that the codimension of its image in $X$ is $\geqslant 2$ and $Z\vert_{X_0}$ is smooth (though possibly disconnected).
Note that our constructions are compatible with restrictions to open subsets. Moreover, since $\jepPubScript{M}(h^{-\alpha})$ is clearly torsion-free, it follows that
$F_k:=F_k\jepPubScript{M}(h^{-\alpha})$ is torsion free, hence the canonical map $F_k\to \iota_*\big(F_k\vert_{X_0}\big)$ is injective. Therefore it is enough to prove the assertion on 
$X_0$, hence we may assume that $Z$ is smooth. However, in this case the assertion follows 
from Corollary~\ref{smooth_case}.
\end{proof}


We can now  define the Hodge ideals for $\jepPubBold{Q}$-divisors. Let $X$ be a smooth complex algebraic variety and $Z$ a reduced effective divisor on $X$.
Given an effective $\jepPubBold{Q}$-divisor $D$ with ${\rm Supp}(D)=Z$, we define coherent ideals sheaves $I_k(D)$  in $\jepPubScript{O}_X$ as follows.
Suppose first that there is a nonzero $h\in\jepPubScript{O}_X(X)$, with $H = {\rm div}(h)$, and a positive rational number $\alpha$ such that $D=\alpha H$. 
It turns out to be more convenient to work with the $\jepPubScript{D}_X$-module $\jepPubScript{M} (h^\beta)$, where $\beta = 1 - \alpha$. Recall that we have a 
filtered isomorphism 
$$\jepPubScript{M} (h^{-\alpha}) \longrightarrow \jepPubScript{M} (h^\beta), \,\,\,\,\,\,wh^{- \alpha} \longmapsto (wh^{-1})h^\beta,$$
and therefore, if $k\geqslant 0$,
it follows from Proposition~\ref{prop_for_def} that there is a unique coherent ideal $I_k(D)$ such that 
$$F_k\jepPubScript{M}(h^{\beta})=I_k(D)\otimes_{\jepPubScript{O}_X} \jepPubScript{O}_X\big(kZ + H \big)h^{\beta}$$
(note that we always have $\lceil D \rceil \geqslant Z$).
The definition is independent of the choice of $\alpha$ and $h$: indeed, using Remark~\ref{rem_behavior_etale}, it is enough to check this after the pullback by a suitable \'{e}tale surjective map, hence we deduce the independence assertion using Remark~\ref{rem_indep_equations1}. 
This implies that the general case of the definition follows by covering $X$ with suitable affine open subsets such that $D$ can be written as above in each of them. Note that when $D=Z$ we have $\beta = 0$, and so the ideals $I_k(D)$ are the Hodge ideals studied in \cite{MP1}.

\begin{remark}\label{rmk_inclusion_ideals}
From the definition and the filtration property,  it follows that we always have the inclusion 
$$\jepPubScript{O}_X (-Z) \cdot I_{k-1}(D) \subseteq I_k (D)\quad\text{for}\quad k\geqslant 1.$$ 
We note that for the reduced divisor $Z$,
we have the more subtle inclusions
$$I_k (Z) \subseteq I_{k-1} (Z)\quad \text{for}\quad k\geqslant 1$$
(see \cite[Prop.~13.1]{MP1}). We do not know however whether this holds for arbitrary $\jepPubBold{Q}$-divisors $D$, 
and in fact we suspect that this is not the case. (Note that it does hold when $D$ has simple normal crossings support
by Proposition \ref{Hodge_ideals_SNC}. It is also shown to hold when $D$ has an isolated weighted homogeneous singularity 
in the upcoming \cite{Zhang}.)
However, when $D = \alpha Z$ these inclusions do hold modulo the ideal $\jepPubScript{O}_X(-Z)$, see \cite[Cor.~B]{MP3}. 
More precisely, we have 
$$I_k(\alpha Z)+\jepPubScript{O}_X(-Z)\subseteq I_{k-1}(\alpha Z)+\jepPubScript{O}_X(-Z)\quad\text{for}\quad k\geqslant 1.$$
This implies in particular that if $I_k(\alpha Z)=\jepPubScript{O}_X$ for some $k\geqslant 1$, then $I_{k-1}(\alpha Z)=\jepPubScript{O}_X$.
\end{remark}

\begin{remark}\label{old_ideal}
According to Proposition \ref{prop_for_def}, we also have ideals $I_k^\prime (D)$ given by
$$F_k\jepPubScript{M}(h^{-\alpha})=I_k^\prime (D)\otimes_{\jepPubScript{O}_X} \jepPubScript{O}_X\big((k+1)Z -\lceil D\rceil \big)h^{-\alpha},$$
which are related to $I_k(D)$ by the formula
$$I_k (D) = I_k^\prime (D) \otimes_{\jepPubScript{O}_X} \jepPubScript{O}_X (Z -\lceil D\rceil).$$
\end{remark}

The following periodicity property often allows us to reduce our study to the case  $\lceil D\rceil = Z$.

\begin{lemma}\label{periodicity}
If $D'$ is an integral divisor with ${\rm Supp}(D')\subseteq {\rm Supp}(D)$, then
$$I_k(D+D')=I_k(D)\otimes_{\jepPubScript{O}_X} \jepPubScript{O}_X (- D').$$
In particular 
$$I_k (D) = I_k (B) \otimes \jepPubScript{O}_X (Z - \lceil D \rceil),$$
with $B = D + Z - \lceil D \rceil$ satisfying $\lceil B \rceil = Z$.
\end{lemma}
\begin{proof}
Using the notation in Remark \ref{old_ideal}, the equivalent statement
$$I_k^\prime (D+D')=I_k^\prime (D)$$
follows from the definition and Remark~\ref{rem_isom_comp_with_filtrations}.
\end{proof}

\begin{remark}\label{always_nontrivial}
Note that $I_k (D) \subseteq \jepPubScript{O}_X (Z - \lceil D \rceil)$ for all $k$, and so if $\lceil D \rceil \neq Z$, then one can never have 
$I_k (D) = \jepPubScript{O}_X$. It is however still interesting to ask whether $I_k (B) = \jepPubScript{O}_X$.
\end{remark}










\section{A global setting for the study of Hodge ideals}\label{global_setting}
We now consider a setting in which we can define global filtered $\jepPubScript{D}_X$-modules that are locally isomorphic to the
$\big(\jepPubScript{M}(h^{-\alpha}), F\big)$ discussed in the previous sections. 

Let $X$ be a smooth variety and $D=(1/\ell)H$ a $\jepPubBold{Q}$-divisor, where $H$ is an integral divisor and $\ell$ is a positive integer.
The extra assumption we make here is that there is a line bundle $M$ such that 
$$\jepPubScript{O}_X(H)\simeq M^{\otimes \ell}.$$ 
We denote by $U$ the complement of $Z={\rm Supp}(H)$
and by $j$ the inclusion $U\hookrightarrow X$.

Let $s\in \Gamma(X,M^{\otimes \ell})$ be a section whose zero locus is $H$. 
Since $s$ does not vanish on $U$, we may consider the section $s^{-1}\in\Gamma\big(U,(M^{-1})^{\otimes \ell}\big)$. Let $p\colon V\to U$ be the \'{e}tale cyclic cover corresponding 
to $s^{-1}$, hence
$$V\simeq {\bf Spec}\big(\jepPubScript{O}_U\oplus M\oplus\cdots\oplus M^{\otimes(\ell-1)}\big).$$
We consider the filtered $\jepPubScript{D}_X$-module $\jepPubScript{M}=j_+p_+\jepPubScript{O}_V$, that underlies
a mixed Hodge module. The  $\mu_{\ell}$-action on $V$, where $\lambda\in\mu_{\ell}$ acts on $M^{\otimes i}$ by multiplication with
$\lambda^{-i}$,
induces an eigenspace decomposition
$$\jepPubScript{M}=\mathop{\textstyle\bigoplus}\limits_{i=0}^{\ell-1}\jepPubScript{M}_i,$$
where $\lambda\in\mu_{\ell}$ acts on $\jepPubScript{M}_i$ by multiplication with $\lambda^{-i}$.
We consider on each $\jepPubScript{M}_i$ the induced filtration.

Note that if $X_0$ is an open subset of $X$ such that we have a trivialization $M\vert_{X_0}\simeq\jepPubScript{O}_{X_0}$, and if 
via the corresponding trivialization of $M^{\otimes \ell}\vert_{X_0}$, the restriction $s\vert_{X_0}$ corresponds to $h_0\in\jepPubScript{O}_X(X_0)$, then 
we have isomorphisms of filtered $\jepPubScript{D}_{X_0}$-modules
$$\jepPubScript{M}_i\simeq\jepPubScript{M}(h_0^{-i/\ell})\quad\text{for $i$ such that}\quad 0\leqslant i\leqslant \ell-1.$$
We also see that the filtration on $\jepPubScript{M}$ is the direct sum filtration, since this holds locally. 
Moreover, we have isomorphisms of $\jepPubScript{O}_{X_0}$-modules
$$\jepPubScript{M}_i\vert_{X_0}\simeq\jepPubScript{O}_X(*Z)\vert_{X_0},$$
which glue to isomorphisms of $\jepPubScript{O}_X$-modules
$$\jepPubScript{M}_i\simeq M^{\otimes i}\otimes_{\jepPubScript{O}_X}\jepPubScript{O}_X(*Z)=j_*j^*M^{\otimes i}.$$
Via these isomorphisms, it follows from the definition of Hodge ideals (see also Remark \ref{old_ideal}) that we have
$$F_k\jepPubScript{M}_i\simeq M^{\otimes i}\otimes_{\jepPubScript{O}_X}I_k^\prime \left(i/\ell \cdot H\right)\otimes_{\jepPubScript{O}_X}\jepPubScript{O}_X\left((k+1)Z-\lceil i/\ell \cdot H\rceil\right)
$$ 
$$\simeq M^{\otimes i} (-H) \otimes_{\jepPubScript{O}_X}I_k\left(i/\ell \cdot H\right)\otimes_{\jepPubScript{O}_X}\jepPubScript{O}_X\left(kZ+ H\right).$$



\section{A complex associated to simple normal crossing divisors}\label{complex_for_SNC}
We now discuss a complex that, as we will see later, gives
a filtered resolution of $\jepPubScript{M}_r(h^{-\alpha})$ by filtered induced $\jepPubScript{D}_X$-modules in the case when $h$ defines a simple normal crossing divisor. 


Let $X$ be a smooth, $n$-dimensional, complex variety, $h\in\jepPubScript{O}_X(X)$ nonzero, and $\alpha$ a nonzero rational number (we allow $\alpha$ 
to be either positive or negative). Let $D = \alpha \cdot {\rm div}(h)$. We denote by $Z$
the support of $D$, and assume that it has simple normal crossings.


Associated to $Z$ we have the following complex of right $\jepPubScript{D}_X$-modules:
$$C^{\jepPubBullet}:\,\,0\longrightarrow \jepPubScript{D}_X\longrightarrow\Omega_X^1(\log Z)\otimes_{\jepPubScript{O}_X}\jepPubScript{D}_X\longrightarrow\cdots\longrightarrow \Omega_X^n(\log Z)\otimes_{\jepPubScript{O}_X}\jepPubScript{D}_X\longrightarrow 0,$$
placed in degrees $-n,\ldots,0$. We denote by $D_i\colon C^i\to C^{i+1}$ its differentials. 
If $x_1,\ldots,x_n$ are local coordinates on $X$, then
$$D_i(\eta\otimes P)=d\eta\otimes P+\sum_{i=1}^n (dx_i\wedge\eta)\otimes \partial_{x_i}P.$$
In fact $C^{\jepPubBullet}$ is a filtered complex, where
$$F_{p-n}C^i=\Omega_X^{i+n}(\log Z)\otimes_{\jepPubScript{O}_X}F_{p+i}\jepPubScript{D}_X.$$
This filtered complex is quasi-isomorphic to the filtered right $\jepPubScript{D}_X$-module $\omega_X(*Z)$ 
corresponding to the filtered left $\jepPubScript{D}_X$-module $\jepPubScript{O}_X(*Z)$ (see \cite[Prop.~3.1]{MP1}, and 
\cite[Prop.~3.11(ii)]{Saito-MHM} for a more general statement).

Given $h$ and $\alpha$ as above, we also consider the filtered complex $C^{\jepPubBullet}_{h^{-\alpha}}$ consisting of the same sheaves, but 
with differential $C^i_{h^{-\alpha}}\to C^{i+1}_{h^{-\alpha}}$ given by
$$D_i-\big((\alpha\cdot {\rm d\log}(h)\wedge \jepPubBullet)\otimes {1_{\jepPubScript{D}_X}}\big).\footnote
{In related settings, for instance involving the de Rham complex of $\jepPubScript{M}(h^{-\alpha})$, this type of complex can already be found in the literature; see for instance \cite[\S6.3.11]{Bjork}.}$$
It is easy to see that this is indeed a filtered complex. 

Suppose now that we also have an effective divisor $T$ supported on $Z$. It is not hard to check that the formula for the map
$$C^i_{h^{-\alpha}}\longrightarrow C^{i+1}_{h^{-\alpha}}$$
induces also a map
$$C^i_{h^{-\alpha}}(-T):=\jepPubScript{O}_X(-T)\otimes_{\jepPubScript{O}_X}\Omega^{i+n}_X(\log Z)\otimes_{\jepPubScript{O}_X}\jepPubScript{D}_X$$
$$\longrightarrow 
C^{i+1}_{h^{-\alpha}}(-T):=\jepPubScript{O}_X(-T)\otimes_{\jepPubScript{O}_X}\Omega^{i+1+n}_X(\log Z)\otimes_{\jepPubScript{O}_X}\jepPubScript{D}_X.$$
This is due to the fact that if locally $T={\rm div}(u)$ and $\eta$ is a local section of $\Omega^{i+n}_X(\log Z)$,
then we can write 
$d(u\eta)=ud(\eta)+u\cdot {\rm d\log}(u)\wedge\eta$.
We thus obtain a filtered subcomplex $C^{\jepPubBullet}_{h^{-\alpha}}(-T)$ of $C^{\jepPubBullet}_{h^{-\alpha}}$.
We emphasize that this is \emph{not} obtained by tensoring $C^{\jepPubBullet}_{h^{-\alpha}}$ with $\jepPubScript{O}_X(-T)$.


\begin{proposition}\label{prop_SNC_complex}
If no coefficient of $D - T$ lies in $\jepPubBold{Z}_{<0}$, then
the complex $C^{\jepPubBullet}_{h^{-\alpha}}(-T)$
is filtered quasi-isomorphic to $\big(h^{-\alpha}\omega_X(*Z), G_\jepPubBullet \big)$, where 
$$G_{k-n}h^{-\alpha}\omega_X(*Z)=0\quad\text{if}\quad k<0,$$
$$G_{-n}h^{-\alpha}\omega_X(*Z)=h^{-\alpha}\omega_X(Z- T)$$
$$\text{and}\quad G_{k-n}h^{-\alpha}\omega_X(*Z)=G_{-n}h^{-\alpha}\omega_X(*Z)\cdot F_k\jepPubScript{D}_X\quad\text{if}\quad k>0.$$
\end{proposition}

\begin{proof}
It is immediate to check that the differential induced on ${\rm gr}^F_pC^{\jepPubBullet}_{h^{-\alpha}}(-T)$ does become 
equal to the differential $D_i$ twisted with the identity on $\jepPubScript{O}_X(-T)$, and therefore for every $p$ we have
$${\rm gr}^F_pC^{\jepPubBullet}_{h^{-\alpha}}(-T)=\jepPubScript{O}_X(-T)\otimes_{\jepPubScript{O}_X}{\rm gr}^F_pC^{\jepPubBullet}.$$
In particular, we have 
$$H^iF_pC^{\jepPubBullet}_{h^{-\alpha}}(-T)=0\quad\text{for every}\quad p\in\jepPubBold{Z}\,\,\text{and}\,\, i\in\jepPubBold{Z}\smallsetminus\{0\},$$
by the result in \cite{MP1} quoted above. Consider now the morphism of right $\jepPubScript{D}_X$-modules
$$\varphi\colon C^0_{h^{-\alpha}}(-T) = \omega_X (Z - T)\otimes_{\jepPubScript{O}_X} \jepPubScript{D}_X \longrightarrow 
h^{-\alpha}\omega_X(*Z)$$
given by 
$$\varphi(w\otimes\eta\otimes Q)=(h^{-\alpha}w\eta)Q.$$
We first check that this morphism is surjective. We do this locally, hence we may assume that we have a system of coordinates $x_1,\ldots,x_n$ on $X$ such that $\jepPubScript{O}_X(-Z)$ is generated by $x_1\cdots x_r$ and $\jepPubScript{O}_X(-T)$ by $x_1^{\beta_1}\cdots
x_r^{\beta_r}$. We also write $h=u x_1^{a_1}\cdots x_r^{a_r},$ where $u$ is an everywhere nonvanishing function, 
and define $\alpha_i=\alpha a_i$ and $\gamma_i=\alpha_i-\beta_i$ for all $i$. Note for later use that
$$\alpha\cdot {\rm d\log}(h)=\frac{\alpha}{u}du+\sum_{i=1}^r\alpha_i\frac{dx_i}{x_i}.$$

The surjectivity of $\varphi$ follows from the fact that 
\begin{equation}\label{eq0_prop_SNC_complex}
{\rm Im}(\varphi)=(h^{-\alpha}x_1^{\beta_1}\cdots x_r^{\beta_r}\eta)\cdot\jepPubScript{D}_X=h^{-\alpha}\omega_X(*Z),
\end{equation}
where $$\eta={\rm d\log}(x_1)\wedge\cdots\wedge
{\rm d\log}(x_r)\wedge dx_{r+1}\wedge\cdots\wedge dx_n,$$ and the second equality in (\ref{eq0_prop_SNC_complex})
   is a consequence of the fact that $-\gamma_i-1\not\in \jepPubBold{Z}_{\geqslant 0}$ for all $i$, by assumption.

In order to complete the proof of the proposition it is enough to show
that, for every $k\geqslant 0$, the following sequence is exact:
$$
\jepPubScript{O}_X(-T)\otimes\Omega_X^{n-1}(\log Z)\otimes F_{k-1}\jepPubScript{D}_X\xrightarrow{\;\textstyle\psi_k\;}\jepPubScript{O}_X(-T)\otimes \omega_X(Z)\otimes F_k\jepPubScript{D}_X$$ 
$$\xrightarrow{\;\textstyle\varphi_k\;} G_{k-n} h^{-\alpha}\omega_X(*Z) \longrightarrow 0,
$$
where $\varphi_k$ is the restriction of $\varphi$ to the $(k-n)$-th level of the filtration and $\psi_k$ is the restriction of the  differential of  $C^{\jepPubBullet}_{h^{-\alpha}}(-T)$.
The surjectivity of $\varphi_k$ is an immediate consequence of the surjectivity of $\varphi$ and the definition of the filtration on $h^{\alpha}\omega_X(*Z)$.



Keeping the above notation for the local coordinates on $X$,
it follows from the definition of $\psi_k$ that
$${\rm Im}(\psi_k)=\prod_{j=1}^rx_j^{\beta_j}\otimes \eta \otimes \bigg(\sum_{i=1}^r\Big(x_i\partial_i-\gamma_i-\frac{\partial u}{\partial x_i}\cdot \frac{\alpha x_i}{u}\Big)\cdot F_{k-1}\jepPubScript{D}_X $$
$$ + \sum_{i=r+1}^n
\Big(\partial_i-\frac{\partial u}{\partial x_i}\cdot\frac{\alpha}{u}\Big)\cdot F_{k-1}\jepPubScript{D}_X\bigg)$$
and it is straightforward to see that this is contained in ${\rm Ker}(\varphi_k)$. We now prove by induction on $k$
that if $\varphi_k(x_1^{\beta_1}\cdots x_r^{\beta_r}\otimes\eta\otimes P)=0$ for some $P\in F_k\jepPubScript{D}_X$, then $x_1^{\beta_1}\cdots x_r^{\beta_r}\otimes\eta \otimes P\in {\rm Im}(\psi_k)$. 
Note that the case $k=0$ is trivial.
Let us write $P=\sum_{u,v}c_{u,v}\partial^{u}x^{v}$, where $u$ and $v$ vary over $\jepPubBold{Z}_{\geqslant 0}^n$. 
After subtracting suitable terms from $P$, we may assume that whenever $c_{u,v}\neq 0$, we have $u_i=0$ for $i>r$.
Furthermore, note that if $u_i,v_i>0$ for some $i\leqslant r$, then we can write
$$\partial^{u}x^{v}=\Big(x_i\partial_i-\gamma_i-\frac{\partial u}{\partial x_i}\cdot \frac{\alpha x_i}{u}\Big)A+B,$$
with both $A$ and $B$ of order $\leqslant k-1$. Therefore we may also assume that whenever $c_{u,v}\neq 0$ and $|u|:=\sum_iu_i=k$, we have
\begin{equation}\label{eq_prop_SNC_complex}
u_iv_i=0\quad\text{for $i$ such that}\quad 1\leqslant i\leqslant n.
\end{equation}



Since 
$$(h^{-\alpha}x_1^{\beta_1}\cdots x_r^{\beta_r}\eta)\partial^{u}x^{v}=(\text{non-zero constant}\footnote{Here we use again the fact that
$-\gamma_i-1\not\in\jepPubBold{Z}_{\geqslant 0}$ for all $i$.})\cdot (h^{-\alpha}x_1^{\beta_1}\cdots x_r^{\beta_r}\eta)x^{v-u},$$
and since (\ref{eq_prop_SNC_complex}) implies that for every $(u,v)$ and $(u',v')$ with $|u|=k$ and $c_{u,v},c_{u',v'}\neq 0$
we have $x^{v-u}\neq x^{v'-u'}$, we conclude that in fact $P\in F_{k-1}\jepPubScript{D}_X$, hence we are done by induction.
\end{proof}










\section{The Hodge ideals of simple normal crossing divisors}\label{SNC}
In this section we show that the Hodge ideals of divisors with simple normal crossing support essentially depend only on the support of the divisor, and therefore can be computed as in \cite[\S8]{MP1}.

\begin{proposition}\label{Hodge_ideals_SNC}
Let $X$ be a smooth variety, and $D$ an effective divisor on $X$ with simple normal crossing  support $Z$. Then for all $k$ we have 
$$I_k(D)=I_k(Z)\otimes_{\jepPubScript{O}_X} \jepPubScript{O}_X(Z - \lceil D \rceil).$$
\end{proposition}

\begin{proof}
Equivalently, we need to show that $I_k^\prime (D) = I_k(Z)$ for all $k$. The assertion is local, hence we may assume that we have coordinates $x_1,\ldots,x_n$
on $X$ such that $Z=H_1+\cdots+H_r$, where $H_i$ is defined by $x_i=0$. The morphism $X\to \jepPubBold{C}^r$ given by 
$(x_1,\ldots,x_r)$ is smooth, hence
using Remark~\ref{rem_behavior_etale} we see that it is enough to prove the proposition when 
$X= \operatorname{Spec}\,\jepPubBold{C}[x_1,\ldots,x_n]$ and $D=\sum_{i=1}^n\alpha_iH_i$,
where $H_i= {\rm div}(x_i)$ and $\alpha_i>0$. Let $\ell$ be the smallest positive integer such that all $a_i:=\ell\alpha_i$ are integers. The assertion to be proved is trivial 
when $\ell=1$, hence from now on we assume $\ell\geqslant 2$. 
Consider the Cartesian diagram 
$$
{\def\labelstyle{\textstyle}\xymatrix{V\ar[r]^{j_0}\ar[d]_{p}&W\ar[d]^{g}\\U\ar[r]_{j}&X,}}
$$
where 
$$j_0\colon V=\operatorname{Spec}\,\jepPubBold{C}[x_1^{\pm 1},\ldots,x_n^{\pm 1},y]/(y^{\ell}-x_1^{-a_1}\cdots x_n^{-a_n})$$
$$\longrightarrow W=\operatorname{Spec}\,\jepPubBold{C}[x_1,\ldots,x_n,z]/(z^{\ell}-x_1^{a_1}\cdots x_n^{a_n}),$$
with $j_0^*(z)=y^{-1}$. 
We will make use of some standard facts about cyclic covers with respect to simple normal crossing divisors,
exploiting the toric variety structure on the normalization of $W$. For basic facts regarding toric varieties, 
we refer to 
\cite{Fulton}.



Let $N$ be the lattice $\jepPubBold{Z}^n$ and $M$ its dual. We also consider the lattice 
$$N'=\{(v_1,\ldots,v_{n+1})\in\jepPubBold{Z}^{n+1}\mid a_1v_1+\cdots+a_nv_n=\ell v_{n+1}\}$$
and its dual 
$$M'=\jepPubBold{Z}^{n+1}/\jepPubBold{Z}\cdot (a_1,\ldots,a_n,-\ell).$$
Note that we have an injective lattice map $N'\to N$, with finite cokernel, induced by the projection onto the first $n$ components, and the dual map $M\to M'$
is again injective, with finite cokernel. In fact, we have an isomorphism $M'/M\simeq\jepPubBold{Z}/ \ell\jepPubBold{Z}$ that maps the class of $(u_1,\ldots,u_{n+1})\in M'$
to the class of $u_{n+1}$ in $\jepPubBold{Z}/\ell \jepPubBold{Z}$. 

We thus have an isomorphism $N'_{\jepPubBold{R}}\simeq N_{\jepPubBold{R}}=\jepPubBold{R}^n$. The strongly convex cone $\sigma=\jepPubBold{R}_{\geqslant 0}^n$ in 
$N_{\jepPubBold{R}}=\jepPubBold{R}^n$ gives the toric variety $X=\jepPubBold{C}^n$. As a cone in $N'_{\jepPubBold{R}}$, 
 $\sigma$ gives an affine toric variety $\widetilde{W}$, and the lattice map $N'\to N$ corresponds to a toric map
$\psi\colon \widetilde{W}\to X$. Note that we have a morphism of $\jepPubScript{O}(X)$-algebras 
$\jepPubScript{O}(W)\to \jepPubScript{O}(\widetilde{W})$ that maps $x_i$ to the element of $\jepPubBold{C}[\sigma^{\vee}\cap M']$
corresponding to the class of the $i$-th element of the standard basis of $\jepPubBold{Z}^n$, and $z$ to the class of $(0,\ldots,0,1)$. It is easy to check that
if we denote by $\lfloor \gamma\rfloor$ the largest integer $\leqslant\gamma$, then
\begin{equation}\label{eq_decomp_W}
\jepPubScript{O}(\widetilde{W})=\mathop{\textstyle\bigoplus}\limits_{0\leqslant j\leqslant\ell-1}\jepPubScript{O} (X)x_1^{-\lfloor j\alpha_1\rfloor}\cdots x_n^{-\lfloor j\alpha_n\rfloor}z^j,
\end{equation}
and consequently to deduce that $\jepPubScript{O}(\widetilde{W})$ is integral over $\jepPubScript{O}(W)$. As the coordinate ring of a toric variety, $\jepPubScript{O}(\widetilde{W})$
is normal, hence it is the integral closure of $\jepPubScript{O}(W)$ in its field of fractions.
Moreover, since $\widetilde{W}$ is a toric variety, we may choose a toric resolution of singularities $Y\to \widetilde{W}$, and let $f\colon Y\to X$ be the composition. Since the map
$Y\to W$ is an isomorphism over the complement of $g^{-1} (\sum H_i)$, it follows that there is an open embedding $\iota\colon V\hookrightarrow Y$ such that $f\circ\iota =j\circ p$. The support $E_Y$ of $Y\smallsetminus
\iota(V)$ is the sum of all prime toric divisors on $Y$. 

The action of the torus $T_{M'}={\rm Spec}\,\jepPubBold{C}[M']$ on $\widetilde{W}$ induces an action of the finite group
${\rm Spec}\,\jepPubBold{C}[M'/M]\simeq {\rm Spec}\,\jepPubBold{C}[\jepPubBold{Z}/\ell\jepPubBold{Z}]=\mu_{\ell}$ on 
$\widetilde{W}$. This is the action induced on the normalization $\widetilde{W}$ by the $\mu_{\ell}$-action on $W$ that we discussed in Section~\ref{coverings}. In particular,
the toric resolution $Y\to\widetilde{W}$ is automatically equivariant. Note that in the decomposition (\ref{eq_decomp_W}), an element $\lambda\in\mu_{\ell}$ acts on the 
summand corresponding to $j$ by multiplication with $\lambda^j$.


The equality $f\circ\iota =j\circ p$ implies that we have an isomorphism of filtered $\jepPubScript{D}_X$-modules
$$j_+p_+\jepPubScript{O}_V\simeq f_+\iota_+\jepPubScript{O}_V=f_+\jepPubScript{O}_Y(*E_Y).$$
As usual, in order to compute the push-forward of $\jepPubScript{O}_Y(*E_Y)$, it is more convenient to work with right $\jepPubScript{D}$-modules.
Recall that there is a complex of right $\jepPubScript{D}_Y$-modules
$$A^\jepPubBullet = A_Y^{\jepPubBullet}:\,\,\,0 \longrightarrow \jepPubScript{D}_Y \longrightarrow \Omega_Y^1(\log E_Y) \otimes_{\jepPubScript{O}_Y} \jepPubScript{D}_Y \longrightarrow \cdots  \longrightarrow \omega_Y(E_Y) \otimes_{\jepPubScript{O}_Y} \jepPubScript{D}_Y \longrightarrow 0$$
located in degrees $-n,\ldots,0$, that is filtered quasi-isomorphic to $\omega_Y(*E_Y)$; see the beginning of Section~\ref{complex_for_SNC}. 
Since $Y$ is a toric variety, we have a canonical isomorphism 
$\Omega_Y^1(\log E_Y)\simeq M'\otimes_{\jepPubBold{Z}}\jepPubScript{O}_Y$ (see \cite[\S4.3]{Fulton}). 
We will also consider the corresponding complex on $X$:
$$A_X^{\jepPubBullet}:\,\,\,0\longrightarrow\jepPubScript{D}_X\longrightarrow M\otimes_{\jepPubBold{Z}}\jepPubScript{D}_X\longrightarrow \cdots\longrightarrow \wedge^nM\otimes_{\jepPubBold{Z}}\jepPubScript{D}_X\longrightarrow 0.$$

It follows from the definition that,  forgetting about the filtration, we have
$$f_+\omega_Y(*E_Y)=\jepPubBold{R} f_*(A^{\jepPubBullet}\otimes_{\jepPubScript{D}_Y}\jepPubScript{D}_{Y\to X}).$$
Note that $\jepPubScript{D}_{Y\to X}=f^* \jepPubScript{D}_X$ as $\jepPubScript{O}_Y$-modules, hence the projection formula implies
$$R^if_*(A^{p-n}\otimes_{\jepPubScript{D}_Y}\jepPubScript{D}_{Y\to X})\simeq \wedge^pM'\otimes_{\jepPubBold{Z}}R^if_*\jepPubScript{O}_Y\otimes_{\jepPubScript{O}_X}\jepPubScript{D}_X=0$$
for $i>0$, since $f$ is the composition of a finite map with a toric resolution. Therefore 
$f_+\omega_Y(*E_Y)$ is represented by the complex $B^{\jepPubBullet}$, where
$$B^{p-n}=\wedge^pM'\otimes_{\jepPubBold{Z}}\psi_*\jepPubScript{O}_{\widetilde{W}}\otimes_{\jepPubScript{O}_X}\jepPubScript{D}_X.$$
In order to describe the differential of this complex, it is convenient to use the isomorphism $M_{\jepPubBold{Q}}\simeq M'_{\jepPubBold{Q}}$
and the decomposition (\ref{eq_decomp_W}). With a little care, it follows from the definitions that if we put
$$B^{p-n}_j=\wedge^pM_{\jepPubBold{Q}}\otimes_{\jepPubBold{Q}}\jepPubScript{O}_X\cdot x_1^{-\lfloor j\alpha_1\rfloor}\cdots x_n^{-\lfloor j\alpha_n\rfloor}z^j\otimes
_{\jepPubScript{O}_X}\jepPubScript{D}_X,$$
then $B^{\jepPubBullet}$ decomposes as the direct sum of the subcomplexes $B^{\jepPubBullet}_j$, for $j$ such that $0\leqslant j\leqslant \ell-1$. 
Furthermore, if we identify each $B^{p-n}_j$ in the obvious way with $A^{p-n}_X$, then the differential 
$$\delta_{B_j}^{p-n}\colon \wedge^pM_{\jepPubBold{Q}}\otimes_{\jepPubBold{Q}}\jepPubScript{D}_X\longrightarrow \wedge^{p+1} M_{\jepPubBold{Q}}\otimes_{\jepPubBold{Q}}\jepPubScript{D}_X$$
is given by 
$$\delta_{B_j}^{p-n}=\delta_{A_X}^{p-n}+(w_j \wedge -)\otimes {\rm Id}_{\jepPubScript{D}_X},$$
where $\delta_{A_X}$ is the differential on $A_X^{\jepPubBullet}$ and
$$w_j=(w_{j,1},\ldots,w_{j,n}),\quad\text{with}\quad w_{j,i}=j\alpha_i-\lfloor j\alpha_i\rfloor.$$
It follows from Proposition~\ref{prop_SNC_complex} that we have a morphism
$$B_j^0\longrightarrow \jepPubScript{M}_r(x_1^{w_{j,1}}\cdots x_n^{w_{j,n}})$$ 
that induces a quasi-isomorphism
$$B_j^{\jepPubBullet}\longrightarrow \jepPubScript{M}_r(x_1^{w_{j,1}}\cdots x_n^{w_{j,n}})$$
(see also Remark~\ref{twist_individual}).

We now bring the filtrations into the picture. It follows from Saito's strictness results (see the discussion in 
Section~\ref{Hodge_modules}; cf. also \cite[\S4,~\S6]{MP1}) that 
$$F_kf_+\omega_Y(*E_Y)={\rm Im}\big(\jepPubBold{R} f_*F_k(A^{\jepPubBullet}\otimes_{\jepPubScript{D}_Y}\jepPubScript{D}_{Y\to X})\longrightarrow \jepPubBold{R} f_*(A^{\jepPubBullet}\otimes_{\jepPubScript{D}_Y}\jepPubScript{D}_{Y\to X})\big).$$
Arguing as above, we deduce that 
$$F_kf_+\omega_Y(*E_Y)={\rm Im}\big(f_*F_k(A^{\jepPubBullet}\otimes_{\jepPubScript{D}_Y}\jepPubScript{D}_{Y\to X})\longrightarrow f_*(A^{\jepPubBullet}\otimes_{\jepPubScript{D}_Y}\jepPubScript{D}_{Y\to X})\big).$$
In other words, $(f_+\omega_Y(*E_Y), F)$ is represented by the filtered complex $B^{\jepPubBullet}$, and using Proposition~\ref{prop_SNC_complex},
we conclude that 
$$f_+\omega_Y(*E_Y)\simeq \mathop{\textstyle\bigoplus}\limits_{j=0}^{\ell-1}\jepPubScript{M}_r(x_1^{w_{j,1}}\cdots x_n^{w_{j,n}}),$$
where the filtration on $\jepPubScript{M}_r(x_1^{w_{j,1}}\cdots x_n^{w_{j,n}})$ is given by 
$$F_{-n}\jepPubScript{M}_r(x_1^{w_{j,1}}\cdots x_n^{w_{j,n}})=x_1^{w_{j,1}}\cdots x_n^{w_{j,n}}\omega_X(Z)\quad\text{and}$$
$$F_{k-n}\jepPubScript{M}_r(x_1^{w_{j,1}}\cdots x_n^{w_{j,n}})=F_{-n}\jepPubScript{M}_r(x_1^{w_{j,1}}\cdots x_n^{w_{j,n}})\cdot F_k\jepPubScript{D}_X\quad\text{for}\quad k\geqslant 1.$$

By comparing the $\mu_{\ell}$-actions, we conclude that the summand 
$\jepPubScript{M}_r(x_1^{-\alpha_1}\cdots x_n^{-\alpha_n})$ on which an element $\lambda\in\mu_{\ell}$ acts by multiplication with $\lambda^{-1}$
corresponds to $j=\ell-1$. Therefore the filtration on 
$\jepPubScript{M}(x_1^{-\alpha_1}\cdots x_n^{-\alpha_n})$ is given by
$$F_{-n}\jepPubScript{M}_r(x_1^{-\alpha_1}\cdots x_n^{-\alpha_n})=(x_1^{-\alpha_1}\cdots x_n^{-\alpha_n})x_1^{\lceil\alpha_1\rceil}\cdots x_n^{\lceil\alpha_n\rceil}\omega_X(Z)\quad\text{and}$$
$$F_{k-n}\jepPubScript{M}_r(x_1^{-\alpha_1}\cdots x_n^{-\alpha_n})=F_{-n}\jepPubScript{M}_r(x_1^{-\alpha_1}\cdots x_n^{-\alpha_n})\cdot F_k\jepPubScript{D}_X\quad\text{for}\quad k\geqslant 1.$$
It is now a straightforward computation to see that
$I_k^\prime(D)$ is the ideal generated by the monomials $\prod_{i=1}^nx_i^{c_i}$, where $0\leqslant c_i\leqslant k$ for all $i$ and $\sum_ic_i=(n-1)k$. 
This coincides with $I_k (Z)$ according to \cite[Prop.~8.2]{MP1},
completing the proof of the proposition.
\end{proof}







\section{Computation in terms of a log resolution}\label{log_resolution}
We use the results of the previous two sections in order to describe Hodge ideals of $\jepPubBold{Q}$-divisors in terms of log resolutions.
Let $X$ be a smooth variety, $h\in\jepPubScript{O}_X(X)$ a nonzero function, $H = {\rm div}(h)$, and $\alpha\in\jepPubBold{Q}_{>0}$. We are interested in computing $I_k(D)$, where $D=\alpha H$.
As always, let $Z={\rm Supp}(D)$ and $\beta = 1 - \alpha$.

Let $f\colon Y\to X$ be a log resolution of the pair $(X, D)$ that is an isomorphism over $U=X\smallsetminus Z$,  
and denote $g=h\circ f\in\jepPubScript{O}_Y(Y)$. 
We fix a positive integer $\ell$ such that $\ell\alpha\in\jepPubBold{Z}$. 
As usual, we consider 
$$p\colon V={\bf Spec}\,\jepPubScript{O}_U[y]/(y^{\ell}-h^{-\ell\alpha})\longrightarrow U$$ 
and the inclusion $j\colon U\hookrightarrow X$. By assumption, we also have an open immersion
$\iota\colon U\hookrightarrow Y$ such that $f\circ\iota=j$. 
 By considering the decompositions of 
$$j_+p_+\jepPubScript{O}_V\simeq f_+\iota_+p_+\jepPubScript{O}_V$$
into isotypical components, we conclude that we have a filtered isomorphism
\begin{equation}\label{birational}
\jepPubScript{M}(h^{-\alpha})\simeq f_+\jepPubScript{M}(g^{-\alpha}).
\end{equation}

We now denote $G=f^*D$, and consider on $Y$ the complex introduced in Section~\ref{complex_for_SNC}:
$$C^{\jepPubBullet}_{g^{-\alpha}}(-\lceil G\rceil):\,\,0\longrightarrow\jepPubScript{O}_Y(-\lceil G\rceil)\otimes_{\jepPubScript{O}_Y} \jepPubScript{D}_Y\longrightarrow\jepPubScript{O}_Y(-\lceil G\rceil)\otimes_{\jepPubScript{O}_Y}\Omega^1_Y(\log E)\otimes_{\jepPubScript{O}_Y}\jepPubScript{D}_Y$$
$$\longrightarrow\cdots\longrightarrow\jepPubScript{O}_Y(-\lceil G\rceil)\otimes_{\jepPubScript{O}_Y}\omega_Y(E)\otimes_{\jepPubScript{O}_Y}\jepPubScript{D}_Y\longrightarrow 0,$$
where $E = (f^*D)_{\rm red}$.
This is placed in degrees $-n,\ldots,0$, and if $x_1,\ldots,x_n$ are local coordinates on $Y$, then its differential is given by
$$\eta\otimes Q\longmapsto d\eta\otimes Q+\sum_{i=1}^n (dx_i\wedge \eta)\otimes\partial_iQ-\big(\alpha\cdot {\rm d\log}(g)\wedge\eta\big)\otimes Q.$$





\begin{theorem}\label{formula_log_resolution}
With the above notation, the following hold:
\begin{enumerate}
\item[(i)] For every $p\neq 0$ and every $k\in \jepPubBold{Z}$, we have
$$R^pf_*\big(C^{\jepPubBullet}_{g^{-\alpha}}(-\lceil G\rceil)\otimes_{\jepPubScript{D}_Y}\jepPubScript{D}_{Y\to X}\big)=0$$
and  
$$R^pf_*F_k\big(C^{\jepPubBullet}_{g^{-\alpha}}(-\lceil G\rceil)
\otimes_{\jepPubScript{D}_Y}\jepPubScript{D}_{Y\to X}\big)=0.$$
\item[(ii)] For every $k\in\jepPubBold{Z}$, the natural inclusion induces an injective map
$$ R^0f_*F_k\big(C^{\jepPubBullet}_{g^{-\alpha}}(-\lceil G\rceil)\otimes_{\jepPubScript{D}_Y}\jepPubScript{D}_{Y\to X}\big)\lhook\joinrel\longrightarrow R^0f_*\big(C^{\jepPubBullet}_{g^{-\alpha}}(-\lceil G\rceil)\otimes_{\jepPubScript{D}_Y}\jepPubScript{D}_{Y\to X}\big).$$
\item[(iii)] We have a canonical isomorphism
$$R^0f_*\big(C^{\jepPubBullet}_{g^{-\alpha}}(-\lceil G\rceil)\otimes_{\jepPubScript{D}_Y}\jepPubScript{D}_{Y\to X}\big)\simeq \jepPubScript{M}_r(h^{-\alpha})$$ 
that induces for every $k\in\jepPubBold{Z}$ an isomorphism
$$R^0f_*F_{k-n}\big(C^{\jepPubBullet}_{g^{-\alpha}}(-\lceil G\rceil)\otimes_{\jepPubScript{D}_Y}\jepPubScript{D}_{Y\to X}\big)\simeq h^{-\alpha}\omega_X\big((k+1)Z-\lceil D\rceil\big)\otimes_{\jepPubScript{O}_X}I_k^\prime (D)$$
$$\simeq h^{\beta}\omega_X\big(kZ + H \big)\otimes_{\jepPubScript{O}_X}I_k(D).$$
\end{enumerate}
\end{theorem}

\begin{proof}
It follows from Lemma~\ref{decomp_filtration}, and from the definition of its filtration, that $\jepPubScript{M}_r(g^{-\alpha})$ is 
a direct summand of a right Hodge $\jepPubScript{D}$-module on $Y$. By Saito's strictness of the filtration of (push-forwards of) such $\jepPubScript{D}$-modules, it follows that for all $k,p\in\jepPubBold{Z}$ the canonical map
$$R^pf_*F_k\big(\jepPubScript{M}_r(g^{-\alpha})\overset{\jepPubBold{L}}{\otimes}_{\jepPubScript{D}_Y}\jepPubScript{D}_{Y\to X}\big)\longrightarrow R^pf_*\big(\jepPubScript{M}_r(g^{-\alpha})\overset{\jepPubBold{L}}{\otimes}_{\jepPubScript{D}_Y}\jepPubScript{D}_{Y\to X}\big)$$
is injective, and its image is equal to
$$F_kR^pf_*\big(\jepPubScript{M}_r(g^{-\alpha})\overset{\jepPubBold{L}}{\otimes}_{\jepPubScript{D}_Y}\jepPubScript{D}_{Y\to X}\big)$$
(see the discussion in Section~\ref{Hodge_modules}).

On the other hand, note that if write 
$G=\alpha\cdot {\rm div}(g)=\sum_i\alpha_iE_i$, then $-\lceil\alpha_i\rceil+\alpha_i\not\in\jepPubBold{Z}_{<0}$ for all $i$.
We may thus apply Proposition~\ref{prop_SNC_complex} for the divisor $G$, with $T=\lceil G\rceil$.
Using Proposition~\ref{Hodge_ideals_SNC} as well,
we see that $C^{\jepPubBullet}_{g^{-\alpha}}(-\lceil G\rceil)$ is filtered quasi-isomorphic to
$\jepPubScript{M}_r(g^{-\alpha})$, hence 
$$R^pf_*\big(C^{\jepPubBullet}_{g^{-\alpha}}(-\lceil G\rceil)\otimes_{\jepPubScript{D}_Y}\jepPubScript{D}_{Y\to X}\big)\simeq R^pf_*\big(\jepPubScript{M}_r(g^{-\alpha})\overset{\jepPubBold{L}}{\otimes}_{\jepPubScript{D}_Y}\jepPubScript{D}_{Y\to X}\big)\quad\text{and}$$
$$R^pf_*F_k\big(C^{\jepPubBullet}_{g^{-\alpha}}(-\lceil G\rceil)\otimes_{\jepPubScript{D}_Y}\jepPubScript{D}_{Y\to X}\big)\simeq 
R^pf_*F_k\big(\jepPubScript{M}_r(g^{-\alpha})\overset{\jepPubBold{L}}{\otimes}_{\jepPubScript{D}_Y}\jepPubScript{D}_{Y\to X}\big).$$

Finally, by the definition of push-forward for right $\jepPubScript{D}$-modules we have
$$R^pf_*\big(\jepPubScript{M}_r(g^{-\alpha})\overset{\jepPubBold{L}}{\otimes}_{\jepPubScript{D}_Y}\jepPubScript{D}_{Y\to X}\big)\simeq 
H^pf_+\jepPubScript{M}_r(g^{-\alpha}),$$
and by (\ref{birational}) this is $0$ if $p\neq 0$, and is canonically isomorphic to $\jepPubScript{M}_r (h^{-\alpha})$ if $p=0$. 
The assertions in the proposition follow by combining all these facts.
\end{proof}
 
\begin{remark}[{\bf Local vanishing}]
The statement in Theorem \ref{formula_log_resolution} i) is a generalization of the Local Vanishing theorem 
for multiplier ideals \cite[Th. 9.4.1]{Lazarsfeld}, in view of the calculation in Proposition 9.1 below.
\end{remark}





\begin{thebibliography}{DMST06}
\bibitem[Bjö93]{Bjork} J.-E. Björk -- \emph{Analytic $\mathscr D$-modules and applications}, Mathematics and its Applications, Kluwer Academic Publisher, Dordrecht, 1993.
\bibitem[Dim04]{Dimca} A. Dimca -- \emph{Sheaves in topology}, Universitext, Springer-Verlag, Berlin, New York, 2004.
\bibitem[DMST06]{Dimca_et_al} A. Dimca, Ph. Maisonobe, M. Saito \& T. Torrelli -- ``Multiplier ideals, $V$-filtrations and transversal sections'', \emph{Math. Ann.} \textbf{336} (2006), no.~4, p.~901--924.
\bibitem[Dut18]{Dutta} Y. Dutta -- ``Vanishing for Hodge ideals on toric varieties'', arXiv:1807.11911, to appear in \emph{Math. Nachr.}, 2018.
\bibitem[Ful93]{Fulton} W. Fulton -- \emph{Introduction to toric varieties}, Ann. of Math. Studies, vol.~131, Princeton University Press, Princeton, NJ, 1993.
\bibitem[GKKP11]{GKKP} D. Greb, S. Kebekus, S. J. Kovács \& T. Peternell -- ``Differential forms on log canonical spaces'', \emph{Publ. Math. Inst. Hautes Études Sci.} \textbf{114} (2011), p.~87--169.
\bibitem[HTT08]{HTT} R. Hotta, K. Takeuchi \& T. Tanisaki -- \emph{D-modules, perverse sheaves, and representation theory}, Progress in Math., vol.~236, Birkhäuser, Boston, Basel, Berlin, 2008.
\bibitem[Laz04]{Lazarsfeld} R. Lazarsfeld -- \emph{Positivity in algebraic geometry. II}, Ergeb. Math. Grenzgeb. (3), vol.~49, Springer-Verlag, Berlin, 2004.
\bibitem[MP16]{MP1} M. Mustaţă \& M. Popa -- ``Hodge ideals'', arXiv:1605.08088, to appear in Memoirs of the AMS, 2016.
\bibitem[MP18a]{MP3} \rule{1.2em}{0.1pt}, ``Hodge ideals for $\jepPubBold{Q}$-divisors, $V$-filtration, and minimal exponent'', arXiv:1807.01935, 2018.
\bibitem[MP18b]{MP2} \rule{1.2em}{0.1pt}, ``Restriction, subadditivity, and semicontinuity theorems for Hodge ideals'', \emph{Internat. Math. Res. Notices} (2018), no.~11, p.~3587--3605.
\bibitem[Pop19]{Popa} M. Popa -- ``$\mathscr D$-modules in birational geometry'', in \emph{Proceedings of the ICM (Rio de Janeiro, 2018), Vol.~2}, World Scientific Publishing Co., Inc., River Edge, NJ, 2019, p.~781--806.
\bibitem[Sai88]{Saito-MHP} M. Saito -- ``Modules de Hodge polarisables'', \emph{Publ. RIMS, Kyoto Univ.} \textbf{24} (1988), no.~6, p.~849--995.
\bibitem[Sai90]{Saito-MHM} \rule{1.2em}{0.1pt}, ``Mixed Hodge modules'', \emph{Publ. RIMS, Kyoto Univ.} \textbf{26} (1990), no.~2, p.~221--333.
\bibitem[Sai07]{Saito-LOG} \rule{1.2em}{0.1pt}, ``Direct image of logarithmic complexes and infinitesimal invariants of cycles'', in \emph{Algebraic cycles and motives. Vol.~2}, London Math. Soc. Lecture Note Ser., vol.~344, Cambridge Univ. Press, Cambridge, 2007, p.~304--318.
\bibitem[Sai09]{Saito-HF} \rule{1.2em}{0.1pt}, ``On the Hodge filtration of Hodge modules'', \emph{Moscow Math. J.} \textbf{9} (2009), no.~1, p.~161--191.
\bibitem[Sai16]{Saito-MLCT} \rule{1.2em}{0.1pt}, ``Hodge ideals and microlocal $V$-filtration'', arXiv:1612.08667, 2016.
\bibitem[Sai17]{Saito-YPG} \rule{1.2em}{0.1pt}, ``A young person's guide to mixed Hodge modules'', in \emph{Hodge theory and $L^2$-analysis}, Adv. Lect. Math. (ALM), vol.~39, Int. Press, Somerville, MA, 2017, p.~517--553.
\bibitem[Zha18]{Zhang} M. Zhang -- ``Hodge filtration for $\jepPubBold{Q}$-divisors with quasi-homogeneous isolated singularities'', arXiv:1810.06656, 2018.
\end{thebibliography}

\end{document}
